\documentclass[11pt,letterpaper,reqno]{amsart}
\usepackage{amsmath,amssymb,amsthm,mathtools}
\usepackage{enumitem}
\usepackage{hyperref}
\usepackage{refcheck}

\addtolength{\hoffset}{-1.5cm}\addtolength{\textwidth}{3cm}
\addtolength{\voffset}{-1.5cm}\addtolength{\textheight}{3cm}

\numberwithin{equation}{section}
\newtheorem{theorem}{Theorem}[section]
\newtheorem{lemma}[theorem]{Lemma}
\newtheorem{proposition}[theorem]{Proposition}
\newtheorem{corollary}[theorem]{Corollary}
\newtheorem{conjecture}[theorem]{Conjecture}
\theoremstyle{remark}
\newtheorem{remark}[theorem]{Remark}

\newcommand{\N}{\mathbb{N}}
\newcommand{\Z}{\mathbb{Z}}
\newcommand{\R}{\mathbb{R}}
\newcommand{\Rcal}{\mathcal{R}}
\newcommand{\eps}{\varepsilon}
\newcommand{\1}{\mathbf 1}
\newcommand{\lowerdens}{\underline{\textup{d}}}
\newcommand{\upperdens}{\overline{\textup{d}}}

\DeclareMathOperator{\lcm}{lcm}
\DeclareMathOperator{\supp}{supp}
\DeclareMathOperator{\rad}{rad}

\renewcommand{\leq}{\leqslant}
\renewcommand{\geq}{\geqslant}
\renewcommand{\le}{\leqslant}
\renewcommand{\ge}{\geqslant}

\allowdisplaybreaks[1]

\setlist[enumerate]{label=\textup{(\roman*)},leftmargin=*}

\title{On the first occurrence of an integer as a prefix sum of divisors}

\author[Chae]{Hyunsik Chae}
\address{Hyunsik Chae, Department of Mathematics, UCLA}
\email{hsch@math.ucla.edu}

\author[Fraiture]{Jimmy Fraiture}
\address{Jimmy Fraiture, London}
\email{jimmy.fraiture@gmail.com}

\author[Hou]{Eric Hou\textsuperscript{*}}
\address{Eric Hou, Los Gatos}
\email{eric.x.hou@gmail.com}

\author[Kova\v{c}]{Vjekoslav Kova\v{c}}
\address{Vjekoslav Kova\v{c}, University of Zagreb Faculty of Science, Department of Mathematics, Bijeni\v{c}ka cesta 30, 10000 Zagreb, Croatia}
\email{vjekovac@math.hr}

\author[Kudeba]{Christian Kudeba\textsuperscript{*}}
\thanks{\textsuperscript{*}Authors marked with an asterisk carried out this work while employed by Principia Math Inc.}
\address{Christian Kudeba, Ottawa, Canada}
\email{christian@principia-math.com}


\author[Shakov]{Anton Shakov\textsuperscript{*}}
\address{Anton Shakov, Ottawa, Canada}
\email{anton@principia-math.com}

\author[Vidal]{Danny Vidal\textsuperscript{*}}
\address{Danny Vidal, Ottawa}
\email{danieljjvidal@gmail.com}

\date{\today}

\begin{document}

\begin{abstract}
For a positive integer $N$, let $f(N)$ be the least positive integer whose increasing list of divisors has an initial segment with sum $N$, when such an integer exists. We prove that $f$ is well-defined for all positive integers $N\neq2,5$. Next, answering a question of Erd\H{o}s, we show that $f(N)=o(N)$ fails in a strong sense, whereas $\limsup_{N\to\infty}f(N)/N=\infty$. Arbitrarily small and arbitrarily large ratios $f(N)/N$ each occur on sets of positive lower density. Moreover, we obtain a doubly exponential upper bound for the lower tail and quantitative lower bounds for the upper tail.
\end{abstract}

\maketitle
\tableofcontents

\section{Introduction}\label{sec:introduction}

Write the positive divisors of an integer $n\geq1$ as
\[
 1=d_1(n)<d_2(n)<\cdots<d_{\tau(n)}(n)=n.
\]
Here $\tau(n)$ is the number of positive divisors of $n$, and
$\sigma(n):=\sum_{d\mid n}d$ denotes their sum.
A positive integer $N$ is \emph{represented by a prefix of divisors} (or simply \emph{represented}) if
\[
 N=d_1(n)+\cdots+d_k(n)
 \quad\text{for some }n\in\N\text{ and } 1\leq k\leq\tau(n).
\]
Let $\Rcal$ be the set of represented positive integers, and define $f\colon\Rcal\to\N$ by
\[
 f(N):=\min\{n\in\N : N=d_1(n)+\cdots+d_k(n)
 \text{ for some }1\leq k\leq\tau(n)\}
\]
for every $N\in\Rcal$.

The obvious inequality $f(\sigma(n))\leq n$, together with the
well-known unboundedness of $\sigma(n)/n$, gives
$\liminf_{N\to\infty,\,N\in\Rcal}f(N)/N=0$; see a comment under \cite{OEIS1054}.
Erd\H{o}s asked whether $f(N)=o(N)$ holds as $N\to\infty$, whether this might at least hold for
almost all $N$ (in an unspecified sense), and whether $\limsup f(N)/N=\infty$.
This problem appears in Guy's collection of open problems in number theory \cite[Problem B2]{Guy} and is
catalogued as Problem~1054 on Bloom's website \emph{Erd\H{o}s problems} \cite{Bloom1054}. Its variant is also mentioned in
Erd\H{o}s and Graham's problem book \cite[p.~93]{ErdosGraham}; the convention there omits the divisor $1$, which only results in an immaterial shift of the function $f$.
The computed values of $f$ are recorded in OEIS sequence A167485 \cite{OEIS1054}, where it is also asked whether
\begin{equation}\label{eq:exact-representability}
 \Rcal=\N\setminus\{2,5\}.
\end{equation}

These questions are not easy to answer because the prefix sums of divisors $d_1(n)+\cdots+d_k(n)$ behave quite erratically.
Nevertheless, we successfully rigorously address all of them by providing nontrivial information on the limiting distribution of the ratios $f(N)/N$.
Moreover, we give a computer-assisted proof of the property \eqref{eq:exact-representability}.


\subsection{Quantitative estimates for tails}

Let $\#S$ denote the cardinality of a finite set $S$.
For $S\subseteq\N$, write
\[
 \lowerdens(S):=\liminf_{X\to\infty}\frac{\#(S\cap[1,X])}{X},
 \quad
 \upperdens(S):=\limsup_{X\to\infty}\frac{\#(S\cap[1,X])}{X}
\]
respectively for \emph{lower} and \emph{upper densities} of $S$.
When these densities agree, their common value is denoted by $\textup{d}(S)$.
Throughout, $\log$ denotes the natural logarithm, and logarithmic
expressions are used only where they are defined.

The first result bounds the frequency
of small ratios $f(N)/N$ uniformly in the cutoff.

\begin{theorem}
\label{thm:small-upper}
There is an absolute constant $c>0$ such that, for all sufficiently small
$\delta>0$ and every $X\geq1$,
\[
 \#\{N\leq X:N\in\Rcal,\ f(N)\leq\delta N\}
 \ll X\exp\{-\exp((1/\delta)^c)\}.
\]
In particular, for every fixed $M>0$ this counting function is
$O_M(\delta^MX)$, uniformly in $X$.
\end{theorem}

An immediate asymptotic consequence of Theorem \ref{thm:small-upper} is
\[ \upperdens\big(\{N\in\Rcal : f(N)\leq\delta N\}\big) \ll \exp\{-\exp((1/\delta)^c)\} \]
for all sufficiently small $\delta>0$.
This fact already strongly refutes Erd\H{o}s's intuition that $f(N)$ is typically much smaller than $N$.

Next, arbitrarily small ratios occur on sets of positive lower density.

\begin{theorem}
\label{thm:small-values}
For every $\delta>0$ there is $c_\delta>0$ such that
\[
 \#\{N\leq X:N\in\Rcal,\ f(N)\leq\delta N\}
 \geq c_\delta X
\]
for all sufficiently large $X$.
\end{theorem}

This time, the asymptotic variant of Theorem \ref{thm:small-values} reads
\[ \upperdens\big(\{N\in\Rcal : f(N)\leq\delta N\}\big) > 0 \]
for every $\delta>0$.

Regarding large ratios $f(N)/N$, our estimate decays more slowly than every
negative power of the threshold.

\begin{theorem}
\label{thm:almost-log-tail}
As $T\to\infty$,
\begin{equation}\label{eq:almost-log-tail}
 \lowerdens(\{N\in\Rcal:N\text{ squarefree},\ f(N)>TN\})
 \geq\left(\frac{e^{-\gamma}}2-o(1)\right)
 \frac{\log\log\log T}{(\log T)(\log\log T)}.
\end{equation}
Consequently, $\{N\in\Rcal:f(N)>TN\}$ has positive lower density
for every fixed $T>0$, and
\[
 \limsup_{\substack{N\to\infty\\N\in\Rcal}}\frac{f(N)}N=\infty.
\]
\end{theorem}

Theorem \ref{thm:almost-log-tail} also settles the last part of Erd\H{o}s's question in the affirmative.
For every fixed $T>0$ the set $\{N\in\Rcal:f(N)>TN\}$ has positive lower density, so which is in an even sharper contrast with Erd\H{o}s's conjecture $f(N)=o(N)$. It remains possible that the density of $\{N:f(N)>TN\}$ tends to zero as $T\to\infty$; Section \ref{sec:tightness} shows that this is equivalent to a certain coverage statement for bounded cofactors, and leaves it open.

The final theorem of this type allows the threshold to grow alongside with $N$, as opposed to fixing $\delta$ or $T$ as before.

\begin{theorem}
\label{thm:subexp-growth}
For every fixed $\eta>0$ there is $c_\eta>0$ such that, for all
sufficiently large $X$,
\begin{equation}\label{eq:subexp-growth}
\begin{aligned}
 \#\biggl\{ & \frac X2<N\leq X:N\in\Rcal,\ N\text{ squarefree},\\
 & \frac{f(N)}N>
 \exp\biggl[\left(\frac1{\sqrt2}-\eta\right)
 \sqrt{(\log\log\log N)(\log\log\log\log N)}\biggr]\biggr\}
 \geq c_\eta\frac X{\log\log\log X}.
\end{aligned}
\end{equation}
For every fixed $s>0$ and $\eta>0$, we have
\begin{equation}\label{eq:positive-moment-growth}
 \frac1X\sum_{\substack{N\leq X\\N\in\Rcal}}
 \left(\frac{f(N)}N\right)^s
 \geq\exp\biggl[\left(\frac{s}{\sqrt2}-\eta\right)
 \sqrt{(\log\log\log X)(\log\log\log\log X)}\biggr]
\end{equation}
for all sufficiently large $X$.
\end{theorem}

\subsection{Hypothetical limiting distribution}

The tail estimates from Theorems \ref{thm:small-upper}--\ref{thm:almost-log-tail} suggest studying the empirical probability measures
\begin{equation}\label{eq:empirical-measures}
 \nu_X:=\frac1{R(X)}
 \sum_{\substack{N\leq X\\N\in\Rcal}}\delta_{f(N)/N},
 \quad R(X):=\#(\Rcal\cap[1,X]),\quad X\geq1,
\end{equation}
where $\delta_x$ is unit mass at $x\in(0,\infty)$.
Since we will see that $R(X)=\lfloor X\rfloor-2$ for $X\geq5$,
the normalization by $R(X)$ is asymptotically equivalent to the one by $X$.
Investigation of weak convergence of the measures \eqref{eq:empirical-measures} can be motivated by the classical analogous results for the sequences $\varphi(N)/N$, see \cite{Schoenberg}, or $\sigma(N)/N$, see \cite{Davenport}.
We find it plausible that the limiting distribution also exists in the case of the sequence $f(N)/N$, but this might be out of reach of the currently available techniques. Still, the present paper attempts to provide some partial results and supporting evidence.

For a factor $d$ of a number $n=ed$, the
integer $e$ is called its \emph{cofactor}. Define
\[
 F_e(d):=\sum_{\substack{r\mid ed\\r\leq d}}r,
 \quad G_E:=\bigcup_{1\leq e\leq E}F_e(\N).
\]
We prove that the family of probability measures $\{\nu_X:X\geq1\}$ is tight on $[0,\infty)$ if and
only if
\[
 \lim_{E\to\infty}\lowerdens(G_E)=1,
\]
but this ``coverage assertion'' remains a conjecture.

The distributional conclusions go substantially beyond tail bounds.
There is one fixed $\sigma$-compact Lebesgue-null set
$\mathcal{Z}\subset(0,\infty)$ carrying the finite part of every
compactified subsequential law. Every such law has no finite atoms,
has $0$ in its support, and has infinite positive moments and infinite
logarithmic mean. In particular, any subsequential probability limit on
$[0,\infty)$ is singular continuous. These conclusions follow from an
explicit limiting measure for all representation witnesses and the
singularity of the distribution of $\sigma(n)/n$.

We also identify the remaining difficulty for the problem of weak convergence of \eqref{eq:empirical-measures}. For $t>0$,
let $r_t^*(N)$ count representations $N=F_e(d)$ with $e\geq2$ and
$ed\leq tN$. Its normalized first moment has a continuous limit
$W^*(t)$, and
\[
 \nu_X([0,t])=W^*(t)-\frac1X\sum_{N\leq X}(r_t^*(N)-1)_++o(1).
\]
Thus, existence of the threshold density is equivalent to convergence of
the displayed surplus term. This surplus has lower limit strictly
larger than $1/9$ for every $t\geq2$. It cannot be discarded and it causes complications.

\subsection{Proof strategies}

For small ratios $f(N)/N$, the main inspiration for the proofs is Tao's unpublished comment under \cite{Bloom1054}. It implies a weaker variant of Theorem \ref{thm:small-upper}, with $\delta^2 X$ on the right-hand side, and it already disproves $f(N)=o(N)$. Following his suggestion, we prove an
estimate for high moments of all prefix sums, which eventually gives Theorem \ref{thm:small-upper}. To obtain
positive density, we use sums of proper divisors of a prescribed abundant
family. The proof adapts the regularity decomposition and collision sieve
of Luca and Pomerance \cite{LucaPomerance}, with Pollack's control of
successive aliquot ratios \cite{Pollack}. The fixed abundant factor makes
the represented ratio arbitrarily small, while the second moment estimate
ensures that a positive proportion of the images are distinct.

For large ratios, the cofactor separates the problem into a moment bound
for large $e$ and an arithmetic obstruction for bounded $e$. Typically
$\sigma(n)$ is divisible by a prescribed fixed modulus. If the target is
rough, the complementary divisor sum therefore has to be rough as well.
Its exact forced modulus makes a large divisor of a finite least common
multiple unavoidable. A Rankin estimate for the resulting divisor tail
yields Theorem \ref{thm:almost-log-tail}, while a uniform divisor-sum normality
estimate and a different choice of parameters give
Theorem \ref{thm:subexp-growth}. We use a single cofactor parametrization
for the large value estimates and the witness measure approach.

For the limiting measures, fix a cofactor and a residue class modulo
$\lcm(1,\ldots,e)$. The reciprocal prefix sum is then
$\sigma(n)/n$ minus a fixed rational number. Davenport's theorem in
arithmetic progressions \cite{PollackPalindromes} and the large-cofactor
moment bound give a convergent witness measure. On $(0,\infty)$, every subsequential ratio law is dominated by that witness measure. Erd\H{o}s's singularity
theorem \cite{ErdosDistribution} consequently transfers to every finite
subsequential law. This argument determines regularity of possible laws
without assuming either tightness or uniqueness.


\subsection{Notation}

We use $s(n):=\sigma(n)-n$ for the sum of the proper divisors,
$\varphi$ for Euler's totient function, $\zeta$ for the Riemann zeta
function, and $\gamma$ for Euler's constant. For integers
$a_1,\ldots,a_k$, we write $\lcm(a_1,\ldots,a_k)$ for their least
common multiple and $\gcd(a,b)$ for the greatest common divisor of $a$ and $b$.
We write $v_p(n)$ for the exponent of a prime $p$ in $n$, and
$p^a\Vert n$ for exact divisibility. The largest and least prime factors of $n$
are $P^+(n)$ and $P^-(n)$, with $P^+(1)=1$ and $P^-(1)=\infty$.
An integer is \emph{$y$-smooth} if $P^+(n)\leq y$, and \emph{$y$-rough} if
$P^-(n)>y$. For real $y\geq1$, we use the \emph{primorial} notation
\[
 y\#:=\prod_{\substack{p\leq y\\p\text{ prime}}}p,
\]
where each prime not exceeding $y$ occurs exactly once. The empty product
is $1$; in particular, $y\#=1$ for $1\leq y<2$. For real $y,u\geq1$, set
\[
 \Delta(y):=\frac{\varphi(y\#)}{y\#}
 =\prod_{\substack{p\leq y\\p\text{ prime}}}\left(1-\frac1p\right),
 \qquad
 \Lambda(u):=\lcm(1,\ldots,\lfloor u\rfloor).
\]
Thus $\Delta(y)$ is the natural density of integers coprime to $y\#$,
and $\Lambda(u)$ is the least common multiple of all positive integers
not exceeding $u$.
The symbols $\pi(x)$ and $\pi(x;q,a)$ count primes and primes in a
reduced residue class, respectively, and
$\psi(x):=\sum_{p^a\leq x}\log p$.
We use $u_+:=\max\{u,0\}$ for the positive part of a real number.
All implied constants are absolute unless their dependence is indicated.
The limits involving a fixed threshold always take $X\to\infty$ first.
Finally, for $j\geq0$ we write
\[
 \sigma_j(n):=
 \begin{cases}
 d_1(n)+\cdots+d_{\tau(n)-j}(n),&0\leq j<\tau(n),\\
 0,&j\geq\tau(n)
 \end{cases}
\]
for the \emph{prefix sum} of all but $j$ largest divisors of $n$. We interpret $\sigma_0$ as $\sigma$.


\section{Arithmetic preliminaries}
\label{sec:arithmetic-preliminaries}

We use the cofactor functions and prime-product notation introduced in
Section \ref{sec:introduction}. Define the complementary reciprocal sum
\[
 g_e(n)=\sum_{\substack{r\mid n\\r\geq e}}\frac1r.
\]
Since the final divisor occurs in its own prefix,
\[
 F_e(d)\geq d,\quad F_e(d)\leq X\ \Longrightarrow\ ed\leq eX.
\]
Mertens' theorem and Chebyshev's estimates give
\[
 \Delta(y)\sim\frac{e^{-\gamma}}{\log y},
 \quad \log(y\#)\ll y,
 \quad \log\Lambda(u)=\psi(u)\leq C_\Lambda u;
\]
see, for example, \cite{Tenenbaum}.

Reflecting the divisor list expresses each prefix sum in terms of the complementary reciprocal sum.

\begin{lemma}
For every $e,d\geq1$,
\begin{equation}\label{eq:reflection}
 F_e(d)=ed\,g_e(ed).
\end{equation}
\end{lemma}
\begin{proof}
The involution $r\mapsto ed/r$ exchanges the divisors $r\geq e$ of $ed$
with its divisors at most $d$.
\end{proof}

The omitted divisors force a congruence whose modulus depends on only
finitely much divisibility information. Fix $e\geq2$ and $d\geq1$, put $n=ed$, and let
\[
 D=\{j:1\leq j<e,\ j\mid n\}
\]
be the set of divisors of $n$ below $e$. Put
\[
 M=\lcm(D),\quad C=\sum_{j\in D}\frac Mj,\quad g=\gcd(e,M),\quad
 K_{e,d}=\frac eg\,C .
\]
Since $1\in D$, these are positive integers. In fact $K_{e,d}\geq2$: if $e/g>1$ this is immediate, while $e/g=1$ implies $C\geq M\geq e\geq2$. For $1\leq j<e$ write
$h=\gcd(j,e)$. Then
\[
 j\mid ed\quad\Longleftrightarrow\quad \frac jh\mid\frac eh\,d
 \quad\Longleftrightarrow\quad \frac jh\mid d,
\]
because $j/h$ and $e/h$ are coprime. Thus $D$ depends only on the residue
of $d$ modulo
\[
 \lcm\bigl(\{j/\gcd(j,e):1\leq j<e\}\bigr).
\]
In particular, for fixed $e$ the integer $K_{e,d}$ takes only finitely many
values.
\begin{lemma}\label{lem:fm-modulus}
With the notation above,
\begin{equation}\label{eq:fm-reflection}
 F_e(d)=\sigma(n)-\frac nM\,C ,
\end{equation}
the integer $M/g$ divides $d$, and
\begin{equation}\label{eq:fm-congruence}
 F_e(d)\equiv\sigma(n)\pmod{K_{e,d}} .
\end{equation}
Moreover $K_{e,d}$ is the largest modulus for which
\eqref{eq:fm-congruence} holds for every $d$ with the same set $D$.
\end{lemma}
\begin{proof}
The map $r\mapsto n/r$ is an involution of the divisors of $n$, and $r>d$
exactly when $n/r<e$. Hence
\[
 \sigma(n)-F_e(d)=\sum_{\substack{r\mid n\\ r>d}}r
 =\sum_{j\in D}\frac nj=\frac nM\sum_{j\in D}\frac Mj=\frac nM\,C .
\]
Every $j\in D$ divides $n$, so $M\mid n$ and $n/M$ is an integer. This
proves \eqref{eq:fm-reflection}.
Since $g=\gcd(e,M)$, the integers $M/g$ and $e/g$ are coprime. From
$M\mid ed$ we get $(M/g)\mid(e/g)d$, hence $(M/g)\mid d$. Therefore
\[
 \frac nM=\frac{ed}M=\frac eg\cdot\frac d{M/g}
\]
is $e/g$ times an integer, and by \eqref{eq:fm-reflection} the difference
$\sigma(n)-F_e(d)$ is a multiple of $(e/g)C=K_{e,d}$.
For the last assertion take $d_0=M/g$, so that $ed_0=\lcm(e,M)$. Every
$j\in D$ divides $M$, hence $ed_0$. Conversely $ed_0$ divides $ed$, so
every divisor of $ed_0$ below $e$ lies in $D$. Thus $d_0$ has the same set
$D$, and at $d_0$ the difference $(ed_0/M)C$ equals $(e/g)C$ exactly. Any
modulus valid for every $d$ with this set $D$ divides that difference,
hence divides $K_{e,d}$.
\end{proof}
If $P^-(d)\geq e$, then $D$ is the set of divisors of $e$ below $e$.
Indeed, for $j<e$ the integer $j/\gcd(j,e)$ is below $e\leq P^-(d)$, so it
divides $d$ only if it equals $1$, that is, only if $j\mid e$. In this
case $M\mid e$, $g=M$, and
\[
 K_{e,d}=\sum_{\substack{j\mid e\\ j<e}}\frac ej=\sigma(e)-1 .
\]
So on $e$-rough inputs, $F_e(d)\equiv\sigma(ed)\pmod{\sigma(e)-1}$.

The sum of divisors is divisible by any fixed positive integer for almost all inputs.

\begin{lemma}
\label{lem:fixed-modulus-normality}
For every fixed positive integer $V$,
\[
 \#\{n\leq Y:V\nmid\sigma(n)\}=o_V(Y)
 \quad(Y\to\infty).
\]
\end{lemma}

\begin{proof}
It suffices to prove the assertion when $V=q$ is a prime power. For every
prime $r\equiv-1\pmod q$, the condition $r\Vert n$ implies
\[
 q\mid r+1=\sigma(r)\mid\sigma(n).
\]
Therefore, if $q\nmid\sigma(n)$, none of these primes divides $n$ exactly
once. For any finite set $\mathcal{P}$ of such primes, the Chinese remainder
theorem gives
\[
 \limsup_{Y\to\infty}\frac1Y
 \#\{n\leq Y:q\nmid\sigma(n)\}
 \leq\prod_{r\in\mathcal{P}}
 \left(1-\frac1r+\frac1{r^2}\right).
\]
The reciprocal sum of the primes in the reduced residue class
$-1\pmod q$ diverges. Letting $\mathcal{P}$ increase through this class makes
the displayed product tend to zero. Factoring a general fixed $V$ into its
prime-power divisors and applying the union bound completes the proof.
\end{proof}

This divisibility property forces the range of the full divisor sum to have density zero.

\begin{lemma}
\label{lem:sigma-range-zero}
The set $\sigma(\N)=\{\sigma(n):n\geq1\}$ has asymptotic density zero.
\end{lemma}

\begin{proof}
If $N=\sigma(n)\leq X$, then $n\leq X$. For every fixed integer $q\geq2$,
the values divisible by $q$ number at most $X/q$, while all remaining values
arise from $o_q(X)$ exceptional sources by
Lemma \ref{lem:fixed-modulus-normality}. Hence
\[
 \upperdens(\sigma(\N))\leq\frac1q.
\]
Letting $q\to\infty$ proves the claim.
\end{proof}

For a prime $q$ and a real number $y\geq1$, define
\[
 B_q(y)=\#\{n\leq y:q\nmid\sigma(n)\}.
\]
The following
estimate follows from a mean-value theorem for nonnegative multiplicative
functions and the Siegel--Walfisz theorem.

\begin{lemma}\label{lem:sigma-rate}
There are absolute constants $c_0,c_1>0$ such that, for all sufficiently
large $y$ and every prime
\[
 3\leq q\leq c_0\frac{\log\log y}{\log\log\log y},
\]
one has
\[
 B_q(y)\ll y\exp\!\left(-c_1\frac{\log\log y}q\right).
\]
Also $B_2(y)\ll\sqrt y$.
\end{lemma}

\begin{proof}
For prime powers define the multiplicative function $h_q$ by
\[
 h_q(p^\nu):=
 \begin{cases}
 0,& p\equiv-1\pmod q\text{ and }\nu=1,\\
 1,& \text{otherwise}.
 \end{cases}
\]
If $p\equiv-1\pmod q$ and $v_p(n)=1$, then
$q\mid\sigma(p)\mid\sigma(n)$. Hence
$\1_{q\nmid\sigma(n)}\leq h_q(n)$. Since $0\leq h_q(p^\nu)\leq1$,
the Halberstam--Richert mean-value bound, in the form
\cite[Lemma~2.4]{Pollack}, applies with $\lambda_1=\lambda_2=1$.
It gives, with an absolute implied constant,
\begin{align*}
 B_q(y)
 &\leq\sum_{n\leq y}h_q(n)\\
 &\ll y\exp\bigg(\sum_{p\leq y}\frac{h_q(p)-1}{p}\bigg)\\
 &=y\exp\bigg(-\sum_{\substack{p\leq y\\p\equiv-1\ (q)}}\frac1p\bigg).
\end{align*}

Choose $z=\exp(Cq^2)$ with $C$ sufficiently large. After decreasing
$c_0$ if necessary, the stated range for $q$ ensures that $z\leq y$ for
all sufficiently large $y$. By the Siegel--Walfisz theorem \cite{Tenenbaum}, uniformly for
$t\geq z$,
\[
 \pi(2t;q,-1)-\pi(t;q,-1)\gg\frac{t}{q\log t}.
\]
Summing over dyadic intervals from $z$ to $y$ therefore gives
\[
 \sum_{\substack{z<p\leq y\\p\equiv-1\ (q)}}\frac1p
 \gg \frac1q\log\frac{\log y}{\log z}
 =\frac1q\bigl(\log\log y-2\log q+O(1)\bigr).
\]
In the stated range the last parenthesis is at least
$\tfrac12\log\log y$ for large $y$, proving the first assertion.
Finally, $\sigma(n)$ is odd exactly when $n$ is a square or twice a square,
so $B_2(y)\ll\sqrt y$.
\end{proof}

The next moment estimate extends the argument of Tao \cite{Bloom1054}, which can be viewed as the first moment approach.
Keeping the condition $e\mid n$ in the moment expansion gives the decay $E^{1-k}$ needed below.

\begin{lemma}\label{lem:moment}
For every integer $k\geq2$ there is a constant $C_k$ satisfying
\begin{equation}\label{eq:Ck-growth}
 \log C_k\leq 2^{k+1}\log(2k+2)
\end{equation}
such that, for all $E,Z\geq1$,
\begin{equation}\label{eq:moment}
 \sum_{n\leq Z}\ \sum_{\substack{e\mid n\\e>E}}g_e(n)^k
 \leq C_kZE^{1-k}.
\end{equation}
Consequently, for all real $A,X\geq1$, the number of integers $N\leq X$
for which there are $e,d\in\N$ satisfying $N=F_e(d)$, $e>E$, and
$ed\leq AN$ is at most
\begin{equation}\label{eq:large-count}
 C_kA^{k+1}E^{1-k}X.
\end{equation}
\end{lemma}

\begin{proof}
Expanding the $k$th power and counting common multiples gives
\begin{align*}
 &\sum_{n\leq Z}\ \sum_{\substack{e\mid n\\e>E}}g_e(n)^k\\
 &\quad=\sum_{e>E}\ \sum_{r_1,\ldots,r_k\geq e}
 \frac1{r_1\cdots r_k}
 \left\lfloor\frac{Z}{\lcm(e,r_1,\ldots,r_k)}\right\rfloor\\
 &\quad\leq Z\sum_{e>E}\ \sum_{r_1,\ldots,r_k\geq e}
 \frac1{r_1\cdots r_k\,\lcm(e,r_1,\ldots,r_k)}.
\end{align*}
Set $x_0=e$, $x_i=r_i$, and $R=x_0\cdots x_k$. Since
$r_1\cdots r_k\geq e^k$ and $e>E$,
\[
 \frac{R^{1/(k+1)}}{r_1\cdots r_k}
 \leq e^{1-k}\leq E^{1-k}.
\]
Consequently,
\[
 \frac1{r_1\cdots r_k\,\lcm(e,r_1,\ldots,r_k)}
 \leq E^{1-k}\frac{R^{-1/(k+1)}}{\lcm(x_0,\ldots,x_k)}.
\]
For a prime $p$ and $t\geq1$, let
$S_{p,t}=\{i:v_p(x_i)\geq t\}$, and for every nonempty
$S\subseteq\{0,\ldots,k\}$ put
\[
 y_S:=\prod_p p^{\#\{t\geq1:S_{p,t}=S\}}.
\]
Then
\[
 R=\prod_S y_S^{|S|},\quad \lcm(x_0,\ldots,x_k)=\prod_S y_S.
\]
Dropping the nesting restrictions on the sets $S_{p,t}$ only enlarges the
sum, and hence
\begin{align*}
 \sum_{x_0,\ldots,x_k\geq1}
 \frac{R^{-1/(k+1)}}{\lcm(x_0,\ldots,x_k)}
 &\leq\prod_{\varnothing\ne S\subseteq\{0,\ldots,k\}}
 \sum_{y\geq1}y^{-1-|S|/(k+1)}\\
 &=\prod_{j=1}^{k+1}
 \zeta\!\left(1+\frac{j}{k+1}\right)^{\binom{k+1}{j}}
 =:C_k.
\end{align*}
This proves \eqref{eq:moment}. Since
$\zeta(1+u)\leq1+u^{-1}$ for $u>0$, summing the binomial exponents gives
\eqref{eq:Ck-growth}.

For \eqref{eq:large-count}, put $n=ed$. By \eqref{eq:reflection},
$N=ng_e(n)$. Such a pair therefore satisfies
$n\leq AX$ and $g_e(n)=N/n\geq1/A$. The number of distinct values $N$ is
at most the number of such pairs $(n,e)$, which by Markov's inequality and
\eqref{eq:moment} is at most
\[
 A^k\sum_{n\leq AX}\ \sum_{\substack{e\mid n\\e>E}}g_e(n)^k
 \leq C_kA^{k+1}E^{1-k}X.
\qedhere \]
\end{proof}


\section{Representability}

We first separate the analytic representability input from the finite
computations needed for the exact classification. This distinction is useful
because the large-value estimates require only the former.

The analytic input shows that almost every odd integer is represented, with a quantitative bound for the exceptional set.

\begin{lemma}
\label{lem:analytic-odd-representability}
There is an absolute $c>0$ such that, for $X\geq3$,
\[
 \#\{n\leq X:n\text{ odd},\ n\notin\Rcal\}
 \ll X^{1-c}+\frac{X}{\log X}.
\]
In particular, this exceptional set has size
$o(X/\log\log\log X)$ as $X\to\infty$.
\end{lemma}

\begin{proof}
The theorem of Montgomery and Vaughan \cite{MVGoldbach}
asserts that at most $O(X^{1-c})$ even integers at most $X$ fail to
be sums of two primes. Discard also integers of the form $2p$, of
which there are $O(X/\log X)$. Each remaining even integer has a
representation $p+q$ with distinct primes $p<q$. The divisors of
$pq$ are $1,p,q,pq$, so $1+p+q$ is a divisor-prefix sum. Translating
by $1$ proves the assertion.
\end{proof}

We now prove the exact representability theorem by combining elementary
constructions, finite verification, and analytic estimates. We identify the
finite checks precisely and use a short rational estimate to extend the
verified range.

\begin{theorem}
\label{thm:fraiture-representability}
The represented positive integers are precisely
\[
 \Rcal=\N\setminus\{2,5\}.
\]
In particular, $f(N)$ is defined for every $N\geq6$.
\end{theorem}

\subsection{Computational verification for finitely many integers}

Subset sums of primes in a suitable interval become divisor-prefix sums after adding $1$.

\begin{lemma}
\label{lem:fraiture-prime-window}
Let $M>1$, and suppose that
$M<p_1<\cdots<p_t\leq Y<M^2$ are distinct primes. If
$J\subseteq\{1,\ldots,t\}$ is nonempty, then
\[
 1+\sum_{j\in J}p_j\in\Rcal.
\]
\end{lemma}

\begin{proof}
Put $n=\prod_{j\in J}p_j$. Every composite divisor of $n$ is larger than
$M^2$, whereas every prime divisor of $n$ is at most $Y<M^2$.
Consequently, the positive divisors of $n$ begin with $1$, followed by the
primes $p_j$, $j\in J$, in increasing order. Their initial sum is the
displayed integer.
\end{proof}

Adjoining a prime extends an interval of subset sums when its translate meets or adjoins the original interval.

\begin{lemma}
\label{lem:fraiture-extension}
Let $C\leq U$ be integers. Suppose that the subset sums of a finite set
of distinct primes contain every integer in $[C,U]$, and let $p$ be a new prime satisfying
$p\leq U-C+1$. After adjoining $p$, the subset sums contain every integer
in $[C,U+p]$.
\end{lemma}

\begin{proof}
The old subset sums cover $[C,U]$, and the subset sums using $p$ cover
$[C+p,U+p]$. The hypothesis gives $C+p\leq U+1$, so these two integer
intervals overlap or are consecutive.
\end{proof}

The finite checks and interval-extension argument cover the following three overlapping ranges.

\begin{proposition}
\label{prop:fraiture-finite}
Every positive integer in each of the following intervals is represented:
\begin{align}
 &[6,10\,000\,000],
 \label{eq:fraiture-small}\\
 &[469\,616,273\,803\,744\,799\,154],
 \label{eq:fraiture-first-window}\\
 &[105\,000\,001,10^{27}+10^8].
 \label{eq:fraiture-large-window}
\end{align}
\end{proposition}

\begin{proof}
The three finite integer checks used in this proof are implemented in
our verifier \texttt{verify.rs} at \cite{FraitureVerifier}. Its committed \texttt{output.json} records a successful
run, together with interval endpoints and sample witnesses. The verifier
constructs the witnesses as it runs; the JSON file does not contain all
$51\,000\,001$ large-seed witnesses.

For \eqref{eq:fraiture-small}, take $1\leq B\leq1000$ and a prime $q>B$.
The divisors of $Bq$ consist first of the divisors of $B$ and then of $q$
times those divisors, in the same order. Thus every integer of the form
\[
 \sigma(B)+q\sum_{i=1}^j d_i(B),
 \quad 1\leq j\leq\tau(B),
\]
is represented. A sieve through $10\,000\,000$ marks every integer in
\eqref{eq:fraiture-small} except $7$, which has the separate
representation $7=1+2+4$ from $n=4$.

For \eqref{eq:fraiture-first-window}, take $M=10\,000$ and
$Y=99\,000\,000<M^2$. The first $94$ primes in $(M,Y]$, lying between
$10\,007$ and $10\,883$, have subset sums covering
$[469\,615,480\,503]$. This is checked by the exact recurrence
$S_0=\{0\}$, $S_j=S_{j-1}\cup(S_{j-1}+p_j)$ for $1\leq j\leq94$.
The verifier checks the hypothesis of
Lemma \ref{lem:fraiture-extension} for each successive prime through
$98\,999\,987$. The resulting subset-sum interval is
\[
 [469\,615,273\,803\,744\,799\,153].
\]
Lemma \ref{lem:fraiture-prime-window} then gives
\eqref{eq:fraiture-first-window} after adding $1$.

For \eqref{eq:fraiture-large-window}, let $M=20\,000\,000$.
For every one of the $51\,000\,001$ integers in
\[
 [105\,000\,000,156\,000\,000],
\]
the verifier constructs and checks a sum of four or five distinct primes
in $(M,2M)$. For completeness, the witness search is deterministic.
For an odd target, first remove the fixed prime $20\,000\,003$;
for an even target, remove nothing. Split the resulting even number into
even summands $e_1+e_2$, beginning with the largest even $e_1$ not exceeding
half that number and decreasing $e_1$ by $2$, subject to
$42\,000\,000\leq e_i\leq78\,000\,000$.
For each $e_i$, search in increasing order for a prime $p$ such that
$p<e_i-p$, both $p$ and $e_i-p$ lie in $(M,2M)$, and neither has already
been used. The verifier checks that this search succeeds for every target,
and then checks each completed witness directly.
The next prime after $2M$ is $40\,000\,003$, which is at
most the inclusive width $51\,000\,001$ of the certified interval. Thus
Lemma \ref{lem:fraiture-extension} starts the interval-extension
induction. After a prime $p$ has been adjoined, the new interval width is
at least $2p$, while Bertrand's postulate supplies the next prime below
$2p$; the induction therefore continues for every subsequent prime.

Set $Y=399\,000\,000\,000\,000<M^2$. It remains to show that
\[
 \sum_{2M<p\leq Y}p>10^{27}+10^8.
\]
A single block suffices. Put $a=Y/2=199\,500\,000\,000\,000$.
Dusart's explicit estimates \cite[Theorem~6.9]{Dusart2010} give
\[
 \frac{x}{\log x}<\pi(x)
 \leq\frac{x}{\log x}\left(1+\frac{1.2762}{\log x}\right)
 <1.05\frac{x}{\log x}
 \quad(x\geq a),
\]
because $\log a>32$. Since $\log Y<34$ and $2M<a$, it follows that
\begin{align*}
 \sum_{2M<p\leq Y}p
 &\geq a\bigl(\pi(Y)-\pi(a)\bigr)\\
 &>a\left(\frac{Y}{34}-\frac{1.05a}{32}\right)\\
 &=\frac{17\,599\,173\,046\,875\,000\,000\,000\,000\,000}{17}
 >1.03\cdot10^{27}>10^{27}+10^8.
\end{align*}
The coarse logarithm inequalities require only
$2.71<e<2.72$ and the rational comparisons
$2.72^{32}<a$ and $Y<2.71^{34}$.
The final bound is an inequality between rational numbers, so this bridge
does not depend on floating-point logarithms or on a rounding convention
in the verifier.
The extended subset-sum interval therefore reaches past $10^{27}+10^8$;
adding $1$ and applying Lemma \ref{lem:fraiture-prime-window} proves
\eqref{eq:fraiture-large-window}.
\end{proof}

For reproducibility, we record the commands and expected output of our finite verifier.
Use the following verifier \cite{FraitureVerifier}.
Run
\begin{quote}
\texttt{rustc -O -C overflow-checks=on verify.rs -o /tmp/v}\par
\texttt{/tmp/v all}
\end{quote}
The final JSON object should have \texttt{"check": "all"} and
\texttt{"ok": true}. The integer checks use a sieve and a subset-sum recurrence to verify
interval overlap and the primality, distinctness, location, and sum of
every large-seed witness. Specifically, the first
check marks $9\,999\,994$ of the $9\,999\,995$ integers in
$[6,10^7]$, with $7$ handled separately; the first prime window contains
$5\,705\,894$ primes; and the large-seed check examines
$51\,000\,001$ targets. The verifier's optional numerical bridge and
tail checks are not needed for the rational and analytic proofs given here.

\subsection{An explicit balanced ternary Goldbach problem}

\begin{lemma}\label{lem:fraiture-balanced-goldbach}
Every odd integer $H\geq10^{27}$ is the sum of three distinct odd primes
$p,q,r$ satisfying
\[
 p,q,r>\frac{H}{30\,000\log H}.
\]
\end{lemma}

\begin{proof}
Helfgott's explicit argument
\cite[\S7.4, equations (7.49)--(7.50)]{Helfgott} gives a real
weighted sum
\[
 \sum_{\substack{p+q+r=H\\p,q,r\ \mathrm{odd\ primes}}}
 (\log p)(\log q)(\log r)
 \eta_+(p/x)\eta_+(q/x)\eta_*(r/x)
 \geq0.000422H^2,
\]
where $x=H/(2+9/(196\sqrt{2\pi}))$, for every odd $H\geq10^{27}$.
Here $\eta_+$ and $\eta_*$ are the real smoothing functions used in
that proof. Their individual smoothing
factors satisfy
\[
 \lVert\eta_+\rVert_\infty\leq1.079955,
 \quad
 \lVert\eta_*\rVert_\infty\leq1.414;
\]
see \cite[equations (7.3) and (7.19)]{Helfgott}. Put
\[
 W=(1.079955)^2\cdot1.414,
 \quad z=\frac{H}{30\,000\log H}.
\]
The explicit estimate
$\psi(t)\leq1.03883t$ of
\cite{RosserSchoenfeld} bounds the total absolute weight of triples with
at least one coordinate at most $z$ by
\begin{align*}
 3W\psi(z)\psi(H)\log H
 &\leq\frac{3W(1.03883)^2}{30\,000}H^2\\
 &<0.000178H^2.
\end{align*}
Triples with repeated coordinates have total absolute weight at most
\[
 3WH(\log H)^3<10^{-20}H^2
 \quad(H\geq10^{27}).
\]
For the last inequality, $(\log H)^3/H$ decreases on this range, and
$\log(10^{27})<63$. Since
$0.000422>0.000178+10^{-20}$, some odd-prime triple remains after both
exceptional classes have been discarded.
\end{proof}

The balanced prime representation now proves representability throughout the remaining infinite range.

\begin{proposition}
\label{prop:fraiture-tail}
Every integer
\[
 n\geq10^{27}+10^8
\]
belongs to $\Rcal$.
\end{proposition}

\begin{proof}
First let $n$ be even and put $H=n-1$. Apply
Lemma \ref{lem:fraiture-balanced-goldbach}, and order its primes as
$p<q<r$. Since $H/(\log H)^2$ increases for $H>e^2$ and
\[
 \frac{10^{27}}{(30\,000\log(10^{27}))^2}>1,
\]
we have
\[
 pq>\left(\frac{H}{30\,000\log H}\right)^2>H>r.
\]
Hence the divisors of $pqr$ begin with $1,p,q,r$, and their sum is
$1+p+q+r=n$.

Now let $n$ be odd. By Bertrand's postulate, choose an odd prime $\ell$
such that
\[
 60\,000\log n<\ell<120\,000\log n,
\]
and put $H=n-1-\ell$. Set $T=10^{27}+10^8$. The function
\[
 h(t)=t-1-120\,000\log t-10^{27}
\]
has positive derivative for $t\geq T$, and $h(T)>0$ because
$\log T<64$. Therefore $H\geq10^{27}$; in the same way,
$H>n/2$. Apply Lemma \ref{lem:fraiture-balanced-goldbach} and order
the resulting primes as $p<q<r$. Since $\log H<\log n$,
\[
 p>\frac{H}{30\,000\log H}
 >\frac{n}{60\,000\log n}
 >120\,000\log n>\ell.
\]
For the penultimate inequality, observe that
$n/(\log n)^2$ increases for $n>e^2$ and already exceeds
$7.2\cdot10^9$ at $n=T$. Finally,
\[
 \ell p>
 60\,000\log n\cdot\frac{H}{30\,000\log H}
 >2H>r.
\]
Thus the divisors of $\ell pqr$ begin with $1,\ell,p,q,r$, whose sum is
$1+\ell+p+q+r=n$.
\end{proof}

\begin{proof}[Proof of Theorem \ref{thm:fraiture-representability}]
The three intervals of Proposition \ref{prop:fraiture-finite} overlap,
and the last one overlaps the infinite interval of
Proposition \ref{prop:fraiture-tail} at $10^{27}+10^8$.
Hence every $n\geq6$ is represented. The integers $1,3,4$ are represented
by $m=1,2,3$, respectively. No prefix can equal $2$: its first term is
$1$, and every longer prefix has sum at least $1+2=3$. Similarly, a
two-term prefix equal to $5$ would have to be $1+4$, but $4\mid m$
implies $2\mid m$, contradicting that $4$ is the second-smallest divisor.
Every prefix of length at least three has sum at least $1+2+3=6$.
Therefore neither $2$ nor $5$ is represented.
\end{proof}


The computational dependencies of the representability argument are as follows.

\begin{remark}
The exact exceptional set in Theorem \ref{thm:fraiture-representability}
uses the three finite integer checks in
Proposition \ref{prop:fraiture-finite}. The infinite-tail
Proposition \ref{prop:fraiture-tail} and the analytic
Lemma \ref{lem:analytic-odd-representability} use none of these checks. In particular,
the analytic proofs of the large-value estimates may use the latter
lemma whenever representedness of odd targets is needed.
\end{remark}


\section{Small values of the least representing integer}

\subsection{Upper bounds for the lower tail}

We extend the first-moment argument from Tao's comment on \cite{Bloom1054} to obtain the following higher-moment estimate.

\begin{lemma}\label{lem:kovac-moment}
There is an absolute constant $C>0$ such that, for every integer $q\geq3$
and every real $x\geq1$,
\[
 \sum_{n\leq x}\sum_{j\geq0}
 \left(\frac{\sigma_j(n)}n\right)^q
 \ll \exp(Cq\log\log q)x.
\]
The implicit constant is absolute.
\end{lemma}

\begin{proof}
Let
\[
 1=r_1<r_2<\cdots<r_{\tau(n)}=n
\]
be the divisors of $n$. Divisor reflection gives
\[
 \frac{\sigma_j(n)}n=\sum_{i>j}\frac1{r_i}.
\]
After expanding the $q$th power and summing over $j$, a fixed tuple
$a_1,\ldots,a_q$ of divisors is counted once for each divisor threshold not
exceeding $\min_i a_i$. Hence
\[
 \sum_{j\geq0}\left(\frac{\sigma_j(n)}n\right)^q
 =
 \sum_{\substack{r,a_1,\ldots,a_q\mid n\\
 r\leq\min(a_1,\ldots,a_q)}}
 \frac1{a_1\cdots a_q}.
\]
Summing over $n\leq x$ and bounding the number of multiples of
$\lcm(r,a_1,\ldots,a_q)$ by\linebreak
$x/\lcm(r,a_1,\ldots,a_q)$, it is enough to prove
\begin{equation}\label{eq:k-S}
 S:=\sum_{\substack{r,a_1,\ldots,a_q\geq1\\r\leq\min_i a_i}}
 \frac1{a_1\cdots a_q\lcm(r,a_1,\ldots,a_q)}
 \ll\exp(Cq\log\log q).
\end{equation}
For each $i$, put
\[
 b_i=\frac r{\gcd(r,a_i)},\quad c_i=\frac{a_i}{\gcd(r,a_i)}.
\]
Then $b_i\mid r$, $\gcd(b_i,c_i)=1$, $a_i=rc_i/b_i$, and
\[
 \lcm(r,a_1,\ldots,a_q)=r\lcm(c_1,\ldots,c_q).
\]
The condition $r\leq a_i$ becomes $b_i\leq c_i$. Dropping the coprimality
conditions and taking $\beta=(q-1)/q$, the inequality
\[
 \frac1{c_i}\leq\frac1{b_i^{1-\beta}c_i^\beta}
 \quad(c_i\geq b_i)
\]
applies term by term. Replacing each divisor $b_i\mid r$ by $r/d_i$
and using $\beta q=q-1$ leaves the power $r^{-2}$. Thus
$S\leq S'S''$, where
\[
 S'=\sum_{r\geq1}\frac1{r^2}
 \left(\sum_{d\mid r}\frac1{d^{(q-1)/q}}\right)^q
\]
and
\[
 S''=\sum_{c_1,\ldots,c_q\geq1}
 \frac1{c_1^{(q-1)/q}\cdots c_q^{(q-1)/q}\lcm(c_1,\ldots,c_q)}.
\]
Both sums have Euler products.

For $S'$, write $S'=\prod_pL_p$, where
\[
 L_p=\sum_{\ell\geq0}p^{-2\ell}
 \left(\sum_{0\leq\gamma\leq\ell}
 p^{-((q-1)/q)\gamma}\right)^q.
\]
For $p\leq q$,
\[
 L_p\leq(1-p^{-2})^{-1}(1-p^{-(q-1)/q})^{-q}.
\]
Since
\[
 \sum_{p\leq q}p^{-(q-1)/q}
 \leq q^{1/q}\sum_{p\leq q}\frac1p
 \ll\log\log q,
\]
one obtains
\[
 \prod_{p\leq q}L_p\ll\exp(Cq\log\log q).
\]
For $p>q$, one has $qp^{-(q-1)/q}\ll1$, and the terms with $\ell\geq1$
give $L_p=1+O(p^{-2})$. Thus
\begin{equation}\label{eq:k-Sprime}
 S'\ll\exp(Cq\log\log q).
\end{equation}

For $S''$, write $S''=\prod_pM_p$, where
\[
 M_p=\sum_{\gamma_1,\ldots,\gamma_q\geq0}
 p^{-((q-1)/q)(\gamma_1+\cdots+\gamma_q)-\max_i\gamma_i}.
\]
Put
\[
 A_h=\sum_{0\leq\gamma\leq h}p^{-((q-1)/q)\gamma}.
\]
Grouping by $h=\max_i\gamma_i$ gives
\[
 M_p=\sum_{h\geq0}p^{-h}(A_h^q-A_{h-1}^q),\quad A_{-1}=0.
\]
For $p\leq q$,
\[
 M_p\leq(1-p^{-1})^{-1}(1-p^{-(q-1)/q})^{-q},
\]
so the product over these primes is $\ll\exp(Cq\log\log q)$. For $p>q$,
the mean-value theorem and $qp^{-(q-1)/q}\ll1$ give
\[
 A_h^q-A_{h-1}^q\ll q p^{-h(q-1)/q}.
\]
Consequently,
\[
 M_p-1\ll q p^{-2+1/q},
\quad
 \sum_{p>q}(M_p-1)\ll q\sum_{n>q}n^{-2+1/q}\ll1.
\]
Hence
\begin{equation}\label{eq:k-Ssecond}
 S''\ll\exp(Cq\log\log q).
\end{equation}
Combining \eqref{eq:k-Sprime} and \eqref{eq:k-Ssecond} proves
\eqref{eq:k-S} and the lemma.
\end{proof}

\begin{proof}[Proof of Theorem \ref{thm:small-upper}]
Let
\[
 \mathcal{E}_\delta(X)=
 \{N\leq X:N\in\Rcal,\ f(N)\leq\delta N\}.
\]
If $\delta X<1$, then $\mathcal{E}_\delta(X)$ is empty, so assume
$\delta X\geq1$. If $N\in\mathcal{E}_\delta(X)$ and $n=f(N)$, then
$n\leq\delta X$ and $N=\sigma_j(n)$ for some $j\geq0$ with
$\sigma_j(n)/n\geq1/\delta$. Therefore, for every integer $q\geq3$,
\begin{align*}
 \#\mathcal{E}_\delta(X)
 &\leq\delta^q\sum_{n\leq\delta X}\sum_{j\geq0}
 \left(\frac{\sigma_j(n)}n\right)^q\\
 &\ll\delta^{q+1}\exp(Cq\log\log q)X.
\end{align*}
Choose a sufficiently small absolute $c>0$ and put
\[
 q=\left\lfloor\exp((1/\delta)^c)\right\rfloor.
\]
For all sufficiently small $\delta$, one has
$C\log\log q\leq\tfrac12\log(1/\delta)$, and hence
\[
 \delta^{q+1}\exp(Cq\log\log q)
 \leq\exp\!\left(-\frac q2\log\frac1\delta\right)
 \leq\exp\!\left(-\exp((1/\delta)^{c/2})\right).
\]
Renaming $c/2$ as $c$ proves the first assertion. The fixed-power bound
follows because this double-exponential expression is $O_M(\delta^M)$ for
every fixed $M>0$.
\end{proof}

\subsection{Where are the small ratios?}

The positive-density assertion for every fixed small ratio is compatible
with the absence of such ratios at any practical search range. The
following bound makes this distinction quantitative.

\begin{proposition}
If $n\in\Rcal$ and $f(n)\leq n/10$, then $n>10^{120}$.
\end{proposition}

\begin{proof}
Put $m=f(n)$. Every divisor-prefix sum of $m$ is at most $\sigma(m)$,
so $f(n)\leq n/10$ implies
\[
 \frac{\sigma(m)}m\geq\frac nm\geq10.
\]
If $m\leq5040$, then
\[
 \frac{\sigma(m)}m=\sum_{d\mid m}\frac1d
 \leq\sum_{d=1}^{5040}\frac1d
 \leq1+\log5040<10,
\]
a contradiction. If $n\leq10^{120}$, then $m\leq10^{119}$; hence
Axler's unconditional refinement of Robin's inequality
\cite[Corollary~2]{AxlerRobin} gives
\[
 \frac{\sigma(m)}m
 <(1+3.15367\cdot10^{-7})e^\gamma\log\log m
 \leq9.99745<10,
\]
again a contradiction.
\end{proof}

The obstruction is not vacuous. For
$m=\lcm(1,\ldots,289)$, direct evaluation of the finite Euler product for
$\sigma(m)/m$ gives
\[
 \frac{\sigma(m)}m=10.004735741\ldots>10.
\]
Thus $n=\sigma(m)$ satisfies
\[
 n=3.078152209\ldots\cdot10^{128},
 \quad
 \frac{f(n)}n\leq\frac m{\sigma(m)}
 =0.099952665\ldots<\frac1{10}.
\]


\subsection{A positive-density family of abundant witnesses}

We now prove Theorem \ref{thm:small-values} by adapting the positive-density
range argument of Luca and Pomerance \cite{LucaPomerance}. The main
additional requirement is that the witnesses have arbitrarily large
abundancy. We enforce it by a fixed divisor $D$. A theorem of Pollack
\cite{Pollack} then bounds the abundancy of the intermediate aliquot sums;
this replaces the deficiency condition used in \cite{LucaPomerance}.
After selecting a regular family, we partition it by its small-prime part
and bound the number of pairs with the same aliquot sum. Cauchy--Schwarz
will convert this collision bound into a positive-density image.

If $n\geq2$, then $s(n)$ is the sum of all proper divisors of $n$, so
$s(n)\in\Rcal$ and
\begin{equation}\label{eq:sv-basic}
 f(s(n))\leq n \quad(n\geq2).
\end{equation}


We will construct $\mathcal{A}(X)\subset[1,X]$ such that
\[
 n/\delta<s(n)\leq C_\delta X\quad(n\in\mathcal{A}(X)),
 \quad \#s(\mathcal{A}(X))\gg_\delta X.
\]
By \eqref{eq:sv-basic}, rescaling $X$ then proves the theorem.

It is enough to prove the theorem for $0<\delta\leq1$. Put
$\lambda=1/\delta$, and choose a fixed squarefree integer $D$ such that
\begin{equation}\label{eq:sv-D}
 \frac{\sigma(D)}D>1+\lambda.
\end{equation}
Such a $D$ exists because $\prod_{p\leq z}(1+1/p)$ is unbounded.

For large $X$, let $\mathcal{A}_0(X)$ consist of the integers
\begin{equation}\label{eq:sv-factorization}
 M=pqrk
\end{equation}
for which $p,q,r$ are prime,
\begin{align*}
 &X^{1/120}<k=Dj\leq X^{1/60},
 \quad \frac{\sigma(j)}j\leq4\zeta(2),\\
 &X^{1/15}<r\leq X^{1/12},\\
 &X^{7/20}<q\leq X^{11/30},\\
 &\frac{X}{2qrk}<p\leq\frac{X}{qrk}.
\end{align*}
Set $\ell=rk$ and $n=q\ell$.

\begin{lemma}
There is a constant $C_\delta>0$ such that, for all sufficiently large
$X$,
\[
 \#\mathcal{A}_0(X)\gg_\delta X,
\]
every $M\in\mathcal{A}_0(X)$ has a unique factorization
\eqref{eq:sv-factorization}, and
\begin{equation}\label{eq:sv-two-sided}
 \lambda<\frac{s(t)}t<C_\delta
 \quad(t\in\{k,\ell,n,M\}).
\end{equation}
Moreover, $q>n^{7/9}$.
\end{lemma}

\begin{proof}
For $J\geq1$,
\[
 \sum_{j\leq J}\frac{\sigma(j)}{j^2}
 =\sum_{d\leq J}\frac1{d^2}\sum_{a\leq J/d}\frac1a
 \leq\zeta(2)\sum_{a\leq J}\frac1a.
\]
Consequently, with $J_1=X^{1/120}/D$ and $J_2=X^{1/60}/D$,
\begin{align*}
 \sum_{\substack{J_1<j\leq J_2\\ \sigma(j)/j\leq4\zeta(2)}}\frac1j
 &\geq \sum_{J_1<j\leq J_2}\frac1j
 -\frac1{4\zeta(2)}\sum_{j\leq J_2}\frac{\sigma(j)}{j^2}\\
 &\geq \left(\frac1{240}+o(1)\right)\log X.
\end{align*}
For fixed $q,r,k$, the interval for $p$ is $(T/2,T]$, where
$T=X/(qrk)\geq X^{8/15}$. The prime number theorem, followed by Mertens'
theorem for primes, therefore gives
\[
 \#\mathcal{A}_0(X)
 \gg \frac{X}{\log X}
 \sum_k\frac1k
 \sum_{X^{1/15}<r\leq X^{1/12}}\frac1r
 \sum_{X^{7/20}<q\leq X^{11/30}}\frac1q
 \gg_\delta X.
\]
The factorization is unique because $p>X^{1/2}$, the prime $q$ lies above
$X^{7/20}$, the prime $r$ lies between $X^{1/15}$ and $X^{1/12}$, and
every prime factor of $k$ is at most $X^{1/60}$.

The function $n\mapsto\sigma(n)/n$ is nondecreasing under divisibility.
Since $D\mid t$ for $t\in\{k,\ell,n,M\}$, \eqref{eq:sv-D} therefore
gives the lower bound in \eqref{eq:sv-two-sided}. The inequality
$\sigma(ab)/(ab)\leq(\sigma(a)/a)(\sigma(b)/b)$ also gives
\[
 \frac{\sigma(k)}k\leq
 \frac{\sigma(D)}D\frac{\sigma(j)}j
 \leq4\zeta(2)\frac{\sigma(D)}D,
\]
and adjoining each of $p,q,r$ multiplies $\sigma(t)/t$ by less than $2$.
This gives the upper bound in \eqref{eq:sv-two-sided}. Finally,
$\ell\leq X^{1/10}$ and $q^2>X^{7/10}\geq\ell^7$, whence
$q^9>q^7\ell^7=n^7$.
\end{proof}

\subsection{Arithmetic regularity and smooth-part classes}

For $u$ sufficiently large, put
\[
 y(u)=\frac{\log\log u}{\log\log\log u},
 \quad Y=y(X^{1/120}).
\]
An integer is called \emph{$Y$-smooth} if all of its prime factors are at most
$Y$.

We collect the almost-everywhere regularity properties of divisor sums established by Luca and Pomerance.

\begin{lemma}\label{lem:LP-inputs}
Put $y(n)=\log\log n/\log\log\log n$. Outside a set of integers of asymptotic density
zero, the following hold:
\begin{enumerate}
\item $v_p(\sigma(n))>v_p(n)$ for every prime $p\leq y(n)$;
\item $P^+(\gcd(n,\sigma(n)))\leq y(n)$;
\item $\sigma(n)/\gcd(n,\sigma(n))$ is divisible by every prime
 $p\leq y(n)$;
\item every prime factor of $s(n)/\gcd(n,\sigma(n))$ exceeds $y(n)$.
\end{enumerate}
If in addition $P^+(n)>n^{7/9}$, then, outside another density-zero set,
\[
 y(n)<q_0\leq n^{10/27},\quad q_0\text{ prime}
 \quad\Longrightarrow\quad q_0^2\nmid s(n).
\]
Finally, outside a set of asymptotic density zero,
\[
 \sum_{\substack{a\mid\sigma(n)\\a>(\log\log n)^2}}\frac1a\leq1.
\]
\end{lemma}

\begin{proof}
The first four properties are the conclusions of
\cite[Lemma~2.1]{LucaPomerance}, with the valuation inequality in \textup{(i)}
stated explicitly in the proof of its part~(4). The square-divisibility
assertion is the consequence of \cite[Lemma~2.2]{LucaPomerance} needed here;
we use it only in the displayed range. The reciprocal-divisor estimate
is \cite[Lemma~2.5]{LucaPomerance}.
\end{proof}

We next count integers with a prescribed smooth part and bound the proportion whose smooth part is large.

\begin{lemma}\label{lem:smooth-part-input}
Suppose that $Y=Y(X)\to\infty$ and $Y=o(\log X)$. For every fixed
$c>0$, uniformly for $Y$-smooth
$d\leq Y^c$,
\[
 \#\{n\leq X:d\text{ is the largest $Y$-smooth divisor of }n\}
 \asymp \frac{X}{d\log Y}.
\]
Moreover, given $\eta>0$, $c$ may be chosen sufficiently large that the
integers $n\leq X$ whose largest $Y$-smooth divisor exceeds $Y^c$ account
for at most $\eta X$ integers for all sufficiently large $X$.
\end{lemma}

\begin{proof}
Write $D_Y(n)$ for the largest $Y$-smooth divisor of $n$. For a
$Y$-smooth integer $d$,
$D_Y(n)=d$ if and only if $n=du$ with $\gcd(u,Y\#)=1$. Counting complete
periods modulo $Y\#$ gives
\[
 \#\{n\leq X:D_Y(n)=d\}
 =\frac Xd\Delta(Y)+O(Y\#).
\]
Chebyshev's estimate gives $Y\#=\exp(O(Y))=X^{o(1)}$, while Mertens'
theorem gives $\Delta(Y)\asymp1/\log Y$. This proves the required
uniform estimate. It is the small-$Y$ form of the counting argument
in \cite[(3.1)]{LucaPomerance}.

For the tail, sum the logarithm of the smooth part:
\[
 \sum_{n\leq X}\log D_Y(n)
 =\sum_{p\leq Y}\log p\sum_{a\geq1}
 \left\lfloor\frac X{p^a}\right\rfloor
 \leq X\sum_{p\leq Y}\frac{\log p}{p-1}
 \ll X\log Y.
\]
The last estimate follows by partial summation from Chebyshev's bound.
Markov's inequality now gives
\[
 \#\{n\leq X:D_Y(n)>Y^c\}\ll X/c.
\]
Choosing $c$ sufficiently large in terms of $\eta$ proves the assertion.
\end{proof}

The following lemma selects a large subfamily with the regularity properties
needed in the collision estimate.

\begin{lemma}\label{lem:sv-regular}
There is a subfamily $\mathcal{A}(X)\subseteq\mathcal{A}_0(X)$ with
$\#\mathcal{A}(X)\gg_\delta X$ having the following properties. For
$M=pqrk\in\mathcal{A}(X)$, let $d$ be the largest $Y$-smooth divisor of
$M$. Then, for every $t\in\{k,\ell,n,M\}$,
\begin{enumerate}
\item $d=\gcd(t,\sigma(t))$, and $d$ is the largest $Y$-smooth divisor of
 both $t$ and $s(t)$;
\item $\sigma(t)/d$ is divisible by every prime at most $Y$, while every
 prime factor of $s(t)/d$ exceeds $y(t)$.
\end{enumerate}
In addition,
\[
 y(n)<q_0\leq n^{10/27},\quad q_0\text{ prime}
 \quad\Longrightarrow\quad q_0^2\nmid s(n),
\]
and
\begin{equation}\label{eq:sv-LP25}
 \sum_{\substack{a\mid\sigma(n)\\a>(\log\log X)^2}}\frac1a\leq1.
\end{equation}
There is also a constant $B_\delta>0$ such that
\begin{equation}\label{eq:sv-image-abundancy}
 \frac{\sigma(s(n))}{s(n)}\leq B_\delta.
\end{equation}
\end{lemma}

\begin{proof}
First discard the integers for which $k$ has a prime factor in
$(Y,y(X)]$. Uniformly for $X^{1/120}<k\leq X^{1/60}$,
\[
 y(k)=y(X)+O\!\left(\frac1{\log\log\log X}\right),
\]
so
\[
 \sum_{Y<p\leq y(X)}\frac1p=o(1).
\]
Indeed,
\[
 \sum_{\substack{k\leq X^{1/60}\\
 p\mid k\text{ for some }Y<p\leq y(X)}}\frac1k
 \ll \log X\sum_{Y<p\leq y(X)}\frac1p=o(\log X).
\]
The discarded $k$ therefore account for $o(X)$ members of
$\mathcal{A}_0(X)$.

Apply Lemma \ref{lem:LP-inputs} to each of $k,\ell,n,M$, with its
large-prime square-divisibility conclusion applied to $n$.

Each excluded set has asymptotic density zero. These exclusions remove only
$o(X)$ members of $\mathcal{A}_0(X)$. To justify this assertion, let
$\mathcal{E}\subseteq\N$ satisfy
$\#(\mathcal{E}\cap[1,T])=o(T)$. Partial summation gives
\begin{equation}\label{eq:sv-harmonic-density-zero}
 \sum_{\substack{n\leq T\\n\in\mathcal{E}}}\frac1n=o(\log T).
\end{equation}
The separated factor ranges make the decompositions
$\ell=rk$, $n=q\ell$ and $M=pn$ unique. An upper-bound estimate for the
prime $p$ then shows that the members for which one of $k,\ell,n$ lies in
$\mathcal{E}$ are
\[
 \ll\frac{X}{\log X}
 \sum_{\substack{t\leq X^{O(1)}\\t\in\mathcal{E}}}\frac1t=o(X),
\]
up to bounded reciprocal prime sums; the deletion at the level $M$ is
directly $o(X)$.

For every retained $t$, Lemma \ref{lem:LP-inputs} and its first four
properties identify $\gcd(t,\sigma(t))$ with the largest $y(t)$-smooth
divisor of $t$ and give the corresponding divisibility properties of
$\sigma(t)$ and $s(t)$. Since $Y\leq y(t)\leq y(X)$, the exclusion made in
the first paragraph and the fact that $p,q,r>y(X)$ show that the same integer
$d$ occurs at all four levels and prove \textup{(i)--(ii)}. The hypothesis
$q>n^{7/9}$ gives the displayed square-divisibility exclusion for $s(n)$,
and the final assertion of Lemma \ref{lem:LP-inputs} gives
\eqref{eq:sv-LP25}.

By \cite[Theorem~1.4]{Pollack}, the fixed exceptional set of integers $n$
for which $s(s(n))/s(n)>s(n)/n+1$ has asymptotic density zero.
One further exclusion at the level $n$, justified by
\eqref{eq:sv-harmonic-density-zero}, therefore ensures that
\[
 \frac{s(s(n))}{s(n)}
 \leq \frac{s(n)}n+1.
\]
Together with \eqref{eq:sv-two-sided}, this yields
\eqref{eq:sv-image-abundancy} with $B_\delta=C_\delta+2$.
\end{proof}

Partition $\mathcal{A}(X)$ according to the integer $d$ in
Lemma \ref{lem:sv-regular}. The image sets belonging to distinct classes
are disjoint, because $d$ is also the largest $Y$-smooth divisor of $s(M)$.

\begin{lemma}\label{lem:sv-classes}
There are positive constants $c,c_1,c_2$, depending at most on $\delta$, and a set $\mathcal{D}_X$ of $Y$-smooth
integers $d\leq Y^c$ such that, writing
\[
 \mathcal{A}_d(X)=\{M\in\mathcal{A}(X):d
 \text{ is the largest $Y$-smooth divisor of }M\},
\]
one has, uniformly for $d\in\mathcal{D}_X$,
\begin{equation}\label{eq:sv-class-size}
 c_1\frac{X}{d\log Y}
 \leq\#\mathcal{A}_d(X)
 \leq c_2\frac{X}{d\log Y},
 \quad
 \sum_{d\in\mathcal{D}_X}\frac1d\gg_\delta\log Y.
\end{equation}
The sets $s(\mathcal{A}_d(X))$, $d\in\mathcal{D}_X$, are pairwise disjoint.
\end{lemma}

\begin{proof}
Fix $\eta=\eta(\delta)>0$ with $\#\mathcal{A}(X)\geq\eta X$.
Apply Lemma \ref{lem:smooth-part-input}, choosing $c$ so large that the
classes with $d>Y^c$ contain at most $\eta X/3$ integers. Since
\[
 \sum_{P^+(d)\leq Y}\frac1d
 =\prod_{p\leq Y}\left(1-\frac1p\right)^{-1}
 \asymp\log Y,
\]
we may choose a sufficiently small $a=a(\delta)>0$ and let $\mathcal{D}_X$
consist of those $d\leq Y^c$ for which
\[
 \#\mathcal{A}_d(X)\geq a\frac{X}{d\log Y}.
\]
The classes outside $\mathcal{D}_X$ then contain at most $\eta X/3$ integers.
Thus the classes in $\mathcal{D}_X$ contain $\gg_\delta X$ integers. The
uniform upper estimate above now gives
\[
 \frac{X}{\log Y}\sum_{d\in\mathcal{D}_X}\frac1d
 \gg_\delta X,
\]
which proves \eqref{eq:sv-class-size}. The disjointness assertion follows
from Lemma \ref{lem:sv-regular}.
\end{proof}

\subsection{The collision estimate}

\begin{proposition}\label{prop:sv-second-moment}
For $d\in\mathcal{D}_X$, put
\[
 R_d(u)=\#\{M\in\mathcal{A}_d(X):s(M)=u\}.
\]
Then, uniformly for $d\in\mathcal{D}_X$,
\begin{equation}\label{eq:sv-second-moment}
 \sum_uR_d(u)^2\ll_\delta\frac{X}{d\log Y}.
\end{equation}
\end{proposition}

\begin{proof}
We prove the collision estimate directly, following the organization of
\cite[(3.4), \S\S3.1--3.2]{LucaPomerance}. All estimates needed for the
present family are recorded below. We first note two reciprocal estimates. Since $d$ is the largest $Y$-smooth divisor of every
member of $\mathcal{A}_d(X)$, the corresponding factor $k$ is $d$ times an
integer with no prime factor at most $Y$. The complete-period count in
Lemma \ref{lem:smooth-part-input}, applied uniformly throughout the fixed
power interval for $k$, and partial summation give
\begin{equation}\label{eq:sv-k-reciprocal}
 \sum_{k}\frac1k\ll\frac{\log X}{d\log Y}.
\end{equation}
Here and below a sum over a factor of a member of $\mathcal{A}_d(X)$ is
restricted to the ranges in \eqref{eq:sv-factorization}. The reciprocal
prime sums over the fixed power intervals for $q$ and $r$ are bounded, so
\begin{equation}\label{eq:sv-m-reciprocal}
 \sum_n\frac1n\ll\frac{\log X}{d\log Y}.
\end{equation}

The diagonal contribution to $\sum_uR_d(u)^2$ is
$\#\mathcal{A}_d(X)$ and is bounded by \eqref{eq:sv-class-size}. Consider
an off-diagonal collision, written as $M=pn$, $M'=p'n'$, with $n\ne n'$.
Then
\begin{equation}\label{eq:sv-collision}
 p\,s(n)+\sigma(n)=p'\,s(n')+\sigma(n').
\end{equation}
By Lemma \ref{lem:sv-regular},
$\gcd(n,\sigma(n))=\gcd(n',\sigma(n'))=d$. Put
\[
 \gcd(s(n),s(n'))=dh.
\]
Every prime factor of $h$ exceeds both $y(n)$ and $y(n')$. Equation
\eqref{eq:sv-collision} shows that $dh$ divides
$\sigma(n)-\sigma(n')$. This difference is nonzero. Indeed, equality
would give $p s(n)=p's(n')$; for all sufficiently large $X$,
\[
 \min(p,p')>\tfrac12X^{8/15}>C_\delta X^{7/15}
 \geq\max\{s(n),s(n')\},
\]
which forces $p=p'$, then $s(n)=s(n')$, and finally $n=n'$. Thus
\[
 A_1=\frac{s(n)}{dh},\quad
 A_2=\frac{s(n')}{dh},\quad
 A_3=\frac{|\sigma(n)-\sigma(n')|}{dh}
\]
are positive integers, with $\gcd(A_1,A_2)=1$.

For fixed $n,n'$, elementary linear Diophantine theory shows that there
are fixed integers $p_0,p'_0$ for which every integral solution has the form
\[
 p=p_0+A_2t,\quad p'=p'_0+A_1t \quad(t\in\Z).
\]
All relevant values of $t$ lie in an interval of length at most
\[
 T:=\frac{Xdh}{n s(n')}.
\]
The bounds
\[
 \lambda n'\leq s(n')\leq C_\delta n',
 \quad nn'\leq X^{14/15},
\]
give $T\gg_\delta X^{1/15}$ and
$T\ll_\delta Xdh/(nn')$, so $\log T\asymp\log X$. The affine forms are
primitive whenever a prime pair exists. For example, a common divisor of
$p_0$ and $A_2$ would divide every value of $p=p_0+A_2t$, whereas
$p>A_2$ for all sufficiently large $X$; the second form is identical.
For odd primes the two primitive forms exclude at most two residue classes,
and the prime $2$ either leaves an allowed parity class or there are no
solutions. The standard two-dimensional upper-bound sieve, in the form used in
\cite[(3.7)]{LucaPomerance}, therefore gives
\[
 \ll_\delta
 \frac{Xdh}{nn'(\log X)^2}
 \frac{A_1}{\varphi(A_1)}
 \frac{A_2}{\varphi(A_2)}
 \frac{A_3}{\varphi(A_3)}
\]
possible prime pairs.

If $v\mid s(t)$, then
\begin{equation}\label{eq:sv-totient}
 \frac v{\varphi(v)}
 \leq\zeta(2)\frac{\sigma(v)}v
 \leq\zeta(2)\frac{\sigma(s(t))}{s(t)}
 \leq\zeta(2)B_\delta
 \quad(t\in\{n,n'\}).
\end{equation}
This bounds the factors belonging to $A_1,A_2$, and also gives
$\varphi(h)\gg_\delta h$.

We now reduce the factor belonging to $A_3$. Let $A_{3,1}$ contain its
prime factors at most $(\log\log X)^2$, let $A_{3,2}$ contain those in
$((\log\log X)^2,\log X]$, and let $A_{3,3}$ contain those exceeding
$\log X$. Since $A_3\ll_\delta X^{7/15}$,
\[
 \frac{A_{3,1}}{\varphi(A_{3,1})}\ll\log Y,
 \quad
 \frac{A_{3,3}}{\varphi(A_{3,3})}\ll1.
\]
Let $A'_{3,2}$ be the largest divisor of $A_{3,2}$ coprime to
$\sigma(n)$. The primes removed from $A_{3,2}$ divide $\sigma(n)$ and
exceed $(\log\log X)^2$, so \eqref{eq:sv-LP25} gives
\[
 \frac{A_{3,2}}{\varphi(A_{3,2})}
 \ll\frac{A'_{3,2}}{\varphi(A'_{3,2})}.
\]
Consequently, it remains to prove
\begin{equation}\label{eq:sv-reduced-collision-sum}
 \frac{X\log Y}{(\log X)^2}
 \sum_{n\ne n'}
 \frac{dh}{nn'}
 \frac{A'_{3,2}}{\varphi(A'_{3,2})}
 \ll_\delta\frac{X}{d\log Y},
\end{equation}
where the sum is restricted to pairs $n,n'$ admitting an off-diagonal
collision in \eqref{eq:sv-collision}, and $dh=\gcd(s(n),s(n'))$.
In particular, $dh\mid\sigma(n)-\sigma(n')$.
Notice also that $A_3$ is divisible by every
prime at most $Y$. Indeed, both $\sigma(n)/d$ and $\sigma(n')/d$ have this
property by Lemma \ref{lem:sv-regular}, while every prime factor of $h$
exceeds $Y$.

Suppose first that $h>X^{10/33}$. Write $n=q\ell$ and $n'=q'\ell'$. We claim that
$\gcd(h,s(\ell)\sigma(\ell))=1$. If a prime $q_0$ divided both $h$
and $\sigma(\ell)$, then the identity
$s(n)=q s(\ell)+\sigma(\ell)$ would imply either $q_0=q$ or
$q_0\mid s(\ell)$. In the first case,
$q_0\mid\gcd(n,\sigma(n))=d$; in the second case,
$q_0\mid\ell$ and $q_0\mid\sigma(\ell)$, so again $q_0\mid d$.
Both conclusions contradict the fact that every prime factor of $h$ exceeds
$Y$, whereas $d$ is $Y$-smooth. If instead $q_0\mid h$ and
$q_0\mid s(\ell)$, then the same identity gives
$q_0\mid\sigma(\ell)$, which has just been excluded. Thus both
$s(\ell)$ and $\sigma(\ell)$ are units modulo $h$. It follows that
$q$ is fixed in a reduced residue class modulo $h$, and the same holds for
$q'$. Moreover,
\[
 n\equiv\sigma(n)\equiv\sigma(n')\equiv n'\pmod h.
\]
Eliminating $q,q'$ gives
\begin{equation}\label{eq:sv-large-h-congruence}
 s(\ell')\ell\sigma(\ell)-s(\ell)\ell'\sigma(\ell')
 \equiv0\pmod h.
\end{equation}
The absolute value of the left side is at most
\[
 2C_\delta(C_\delta+1)\max(\ell,\ell')^3
 \ll_\delta X^{3/10}=o(X^{10/33}),
\]
so it vanishes. For every prime power $q_0^a\Vert d$,
Lemma \ref{lem:sv-regular} gives
$q_0^a\Vert s(\ell)$ and $q_0^{a+1}\mid\sigma(\ell)$; it also excludes
primes exceeding $Y$ from $\gcd(\ell,s(\ell))$. Hence
\[
 \gcd(\ell^2,s(\ell))=d,
 \quad
 \gcd((\ell')^2,s(\ell'))=d.
\]
The equality in \eqref{eq:sv-large-h-congruence} can be written as
\[
 \frac{\ell^2}{s(\ell)}+\ell
 =\frac{(\ell')^2}{s(\ell')}+\ell'.
\]
The reduced denominators of the two fractions are $s(\ell)/d$ and
$s(\ell')/d$. Since their difference is an integer, these denominators
are equal. Thus $s(\ell)=s(\ell')$, and the displayed
equality, viewed as a strictly increasing quadratic in $\ell$, gives
$\ell=\ell'$.

In this range Mertens' theorem gives
$A'_{3,2}/\varphi(A'_{3,2})\ll\log\log X/\log Y$. Summing the left side of
\eqref{eq:sv-reduced-collision-sum} and using $\ell=\ell'$ gives
\[
 \ll\frac{Xd\log\log X}{(\log X)^2}
 \sum_{q,q',\ell,h}\frac{h}{qq'\ell^2}.
\]
We may assume $q>q'$ by symmetry. Because $\ell=\ell'$, the two primes
belong to the same residue class modulo $h$. Thus $q=q'+jh$ with $j\geq1$,
and
\[
 \sum_q\frac1q\leq\sum_{1\leq j\leq X}\frac1{jh}
 \ll\frac{\log X}{h},
\]
even if primality is discarded. Summing next
over $h\mid s(q'\ell)$ is bounded by
$\tau(s(q'\ell))=X^{o(1)}$. The reciprocal sum over $q'$ is bounded and
$\sum_\ell\ell^{-2}\ll X^{-1/15}$. The contribution is therefore
\[
 dX^{14/15+o(1)}=o\!\left(\frac{X}{d\log Y}\right),
\]
uniformly for $d\leq Y^c$.

It remains to treat $h\leq X^{10/33}$. Let $A''_{3,2}$ be the largest
divisor of $A'_{3,2}$ coprime to $s(n')$. Equation
\eqref{eq:sv-totient} gives
\[
 \frac{A'_{3,2}}{\varphi(A'_{3,2})}
 \ll_\delta\frac{A''_{3,2}}{\varphi(A''_{3,2})}.
\]
We first note that $h$ is squarefree in this range. Indeed, if
$q_0^2\mid h$, then $q_0^2\mid s(n)$ and $q_0>y(n)$. By
Lemma \ref{lem:LP-inputs}, this forces $q_0>n^{10/27}$. Since
$n>X^{17/40}$,
\[
 q_0^2>n^{20/27}>X^{17/54}>X^{10/33}\geq h,
\]
a contradiction.

Fix $n',h,k$. The congruences
\[
 qrk\equiv(q+1)(r+1)\sigma(k)\equiv n'\pmod h
\]
may therefore be solved for the required symmetric functions. Indeed, any prime
dividing both $h$ and either $n$ or $\sigma(n)$ would divide
$\gcd(n,\sigma(n))=d$, which is impossible because every prime factor of
$h$ exceeds $Y$. Thus $k$ and $\sigma(k)$ are units modulo $h$. The two
congruences therefore
determine $qr$ and $q+r$ modulo $h$. Since $h$ is squarefree, there are at
most $\tau(h)$ possible ordered residue pairs
\[
 q\equiv a\pmod h,\quad r\equiv b\pmod h.
\]
For one such pair put
\[
 \mathfrak{f}(qr)=
 \sum_{\substack{(\log\log X)^2<q_0\leq\log X\\
 q_0\mid\sigma(kqr)-\sigma(n')\\
 q_0\nmid h\sigma(kqr)}}\frac1{q_0},
\]
where $q_0$ is prime. Mertens' theorem gives
\begin{equation}\label{eq:sv-intermediate-totient}
 \frac{A''_{3,2}}{\varphi(A''_{3,2})}
 \ll_\delta 1+\frac{\log\log X}{\log Y}\mathfrak{f}(qr).
\end{equation}
Indeed, when $\mathfrak{f}(qr)\leq1$, the logarithm of the left side is
$O(\mathfrak{f}(qr))$. When $\mathfrak{f}(qr)>1$, Mertens' theorem bounds
the contribution of the full prime range by
$O(\log\log X/\log Y)$.

For fixed $n',h,k,r$ and a prime $q_0$ occurring in
$\mathfrak{f}(qr)$, the congruence
\[
 (q+1)\sigma(kr)\equiv\sigma(n')\pmod{q_0}
\]
fixes $q$ modulo $q_0$. The condition
$q_0\nmid\sigma(kqr)$ ensures that $\sigma(kr)$ is invertible modulo
$q_0$. Since $q_0\nmid h$, this congruence combines with
$q\equiv a\pmod h$. Brun--Titchmarsh and partial summation, followed by
\eqref{eq:sv-totient}, give
\begin{equation}\label{eq:sv-f-sum}
 \sum_{q,r}\frac{\mathfrak{f}(qr)}{qr}
 \ll_\delta
 \begin{cases}
 \dfrac{\log X}{hX^{1/20}},&X^{1/20}<h\leq X^{10/33},\\[6pt]
 \dfrac{1}{h^2(\log\log X)^2},&h\leq X^{1/20}.
 \end{cases}
\end{equation}
Here the sums are over the fixed residue pair. To see
\eqref{eq:sv-f-sum}, first note that a nonreduced residue class modulo
$q_0$ contains no admissible $q$, because $q>X^{7/20}>\log X\geq q_0$.
Otherwise the combined class modulo $q_0 h$ is reduced. Uniformly in
the present ranges,
$q_0 h\leq X^{10/33}\log X<X^{7/20-\epsilon_0}$ for a fixed
$\epsilon_0>0$, so Brun--Titchmarsh and partial summation bound its
reciprocal prime sum by
$O(1/\varphi(q_0 h))\ll_\delta1/(q_0 h)$. Summing the extra weight
$1/q_0$ gives $O(1/(h(\log\log X)^2))$. The reciprocal sum over $r$ in its
class is
\[
 \ll_\delta
 \begin{cases}
 (\log X)X^{-1/20},&h>X^{1/20},\\
 h^{-1},&h\leq X^{1/20},
 \end{cases}
\]
which proves the estimate. The same progression bounds without
$\mathfrak{f}$ and \eqref{eq:sv-intermediate-totient} yield
\begin{equation}\label{eq:sv-qr-sum}
 \sum_{q,r}\left(\frac1{qr}
 +\frac{\log\log X}{\log Y}\frac{\mathfrak{f}(qr)}{qr}\right)
 \ll_\delta
 \begin{cases}
 \dfrac{(\log X)(\log\log X)}{hX^{1/20}\log Y},
 &X^{1/20}<h\leq X^{10/33},\\[6pt]
 h^{-2},&h\leq X^{1/20}.
 \end{cases}
\end{equation}

Using \eqref{eq:sv-k-reciprocal}, the contribution for fixed $n',h$ to
\eqref{eq:sv-reduced-collision-sum} is at most
\[
 \frac{X\log Y}{(\log X)^2}\frac{dh}{n'}
 \frac{\log X}{d\log Y}\tau(h)
 \times\text{the applicable bound in \eqref{eq:sv-qr-sum}}.
\]
For $X^{1/20}<h\leq X^{10/33}$ this is
\[
 \ll_\delta
 \frac{X^{19/20}\log\log X}{\log Y}\frac{\tau(h)}{n'}.
\]
Now $\sum_{h\mid s(n')}\tau(h)\leq\tau(s(n'))^2=X^{o(1)}$, and
\eqref{eq:sv-m-reciprocal} shows that the total is
\[
 \frac{X^{19/20+o(1)}}d
 =o\!\left(\frac{X}{d\log Y}\right).
\]
For $h\leq X^{1/20}$ the fixed-$n',h$ contribution is
\[
 \ll_\delta\frac{X}{\log X}\frac{\tau(h)}{hn'}.
\]
Finally,
\[
 \sum_{h\mid s(n')}\frac{\tau(h)}h
 \leq\left(\frac{\sigma(s(n'))}{s(n')}\right)^2
 \leq B_\delta^2.
\]
Summing over $h$ and then using \eqref{eq:sv-m-reciprocal} gives
$O_\delta(X/(d\log Y))$. Together with the large-$h$ and diagonal
contributions, this proves \eqref{eq:sv-second-moment}.
\end{proof}

\begin{proof}[Proof of Theorem \ref{thm:small-values}]
For each $d\in\mathcal{D}_X$, Cauchy--Schwarz,
Lemma \ref{lem:sv-classes}, and Proposition \ref{prop:sv-second-moment} give
\[
 \#s(\mathcal{A}_d(X))
 \geq\frac{\#\mathcal{A}_d(X)^2}{\sum_uR_d(u)^2}
 \gg_\delta\frac{X}{d\log Y}.
\]
The image sets are disjoint, and hence
\[
 \#s(\mathcal{A}(X))
 \gg_\delta\frac{X}{\log Y}
 \sum_{d\in\mathcal{D}_X}\frac1d
 \gg_\delta X.
\]
If $N=s(M)$ with $M\in\mathcal{A}(X)$, then
\eqref{eq:sv-basic} and \eqref{eq:sv-two-sided} give
\[
 f(N)\leq M<\delta N.
\]
Moreover, $N=s(M)<C_\delta M\leq C_\delta X$. Replacing $X$ by $X/C_\delta$ proves the theorem
for $0<\delta\leq1$, and the case $\delta>1$ follows from $\delta=1$.
\end{proof}


\section{Upper tails}

The moment bound handles representations with large cofactors. For small
cofactors, we use targets that are coprime to all small primes. Divisibility
of $\sigma(n)$ then forces the omitted-divisor kernel to be unusually large.
The following estimate quantifies the scarcity of those kernels.

\subsection{Divisor tails of a least common multiple}

\begin{lemma}
\label{lem:kernel-tails}
For $y,H\geq2$, put
\[
 S(y,H):=\sum_{\substack{L\mid \Lambda(y)\\L>H}}\frac{\tau(L)}L.
\]
If $E\to\infty$ and
\[
 P:=\left\lfloor
 E^{\,1+(2+(\log\log\log E)^{-1/2})\log\log E/\log\log\log E}
 \right\rfloor,
\]
then
\begin{equation}\label{eq:fixed-kernel-tail}
 \log S(E,P/E)
 \leq-\frac{\log\log E}{\sqrt{\log\log\log E}}
 +O\left(\frac{\log\log E}{\log\log\log E}\right).
\end{equation}
If instead $J\to\infty$,
\[
 W=\sqrt{\frac{J\log J}{2}},
 \quad F=\lfloor e^W\rfloor,
\]
then, for every fixed $C>0$,
\begin{equation}\label{eq:moving-kernel-tail}
 S\left(F,\frac{e^J}{CW}\right)
 \leq\exp\{-W+o(W)\}.
\end{equation}
\end{lemma}

\begin{proof}
For $0<\alpha\leq1/4$, Rankin's inequality and the Euler product for the
divisor function give
\begin{align}
 S(y,H)
 &\leq H^{-\alpha}
 \prod_{p\leq y}\sum_{a\geq0}
 \frac{a+1}{p^{a(1-\alpha)}}
 =H^{-\alpha}
 \prod_{p\leq y}(1-p^{-1+\alpha})^{-2}.
 \label{eq:rankin-kernel}
\end{align}
Mertens' estimate for $\sum_{p\leq y}1/p$, Chebyshev's bound, and partial
summation show that, when $z=\alpha\log y\to\infty$,
\begin{align}
 \sum_{p\leq y}p^{-1+\alpha}
 &=\log\log y+
 O\left(\int_0^z\frac{e^u-1}{u}\,du\right)
 =\log\log y+O\left(\frac{e^z}{z}\right).
 \label{eq:sharp-prime-sum}
\end{align}
The higher terms in the logarithm of the Euler product in
\eqref{eq:rankin-kernel} are $O(1)$ uniformly for $\alpha\leq1/4$.

For \eqref{eq:fixed-kernel-tail}, take
$\alpha=\log\log\log E/\log E$. Then
\[
 \alpha\log(P/E)
 =\left(2+\frac1{\sqrt{\log\log\log E}}\right)\log\log E+o(1),
\]
whereas \eqref{eq:sharp-prime-sum} gives
\[
 2\sum_{p\leq E}p^{-1+\alpha}
 =2\log\log E+O\left(\frac{\log\log E}{\log\log\log E}\right).
\]
Substitution in \eqref{eq:rankin-kernel} proves the first assertion.

For \eqref{eq:moving-kernel-tail}, take $\alpha=W/J$. Since
\[
 \alpha\log F=\frac12\log J+o(1),
 \quad
 \alpha\log\left(\frac{e^J}{CW}\right)=W+o(W),
\]
equation \eqref{eq:sharp-prime-sum} gives
\[
 \sum_{p\leq F}p^{-1+\alpha}
 \ll\log W+\frac{\sqrt J}{\log J}=o(W).
\]
The second assertion follows once more from \eqref{eq:rankin-kernel}.
\end{proof}

\subsection{Periodic obstructions for bounded cofactors}

For sufficiently large $A$, abbreviate the upper-tail scale by
\[
 L(A):=\frac{\log\log\log A}{(\log A)(\log\log A)}.
\]
The forced moduli identify a periodic set that almost completely avoids
the bounded-cofactor range. Its density estimate will also be useful for
the coverage questions in Section \ref{sec:tightness}.

For $A\geq2$ let $\mathcal{K}_A$ be the set of moduli $K_{e,d}$ with
$2\leq e\leq A$ and $d\geq1$; it is finite by
Section \ref{sec:arithmetic-preliminaries}. Put
\[
 V_A=\{N\in\N:k\nmid N\ \text{for every }k\in\mathcal{K}_A\} .
\]
This is a union of residue classes modulo $W_A=\lcm(\mathcal{K}_A)$, so it has
a natural density, and $\textup{d}(V_A)\geq\varphi(W_A)/W_A>0$. By inclusion and
exclusion the density is the rational number
\[
 \textup{d}(V_A)=1-\sum_{\varnothing\neq I\subseteq\mathcal{K}_A}
 \frac{(-1)^{|I|+1}}{\lcm(I)} .
\]
\begin{proposition}\label{prop:fm-envelope}
Fix $A\geq2$. All but $o_A(X)$ of the integers $N\leq X$ in $G_A$ are
divisible by some $k\in\mathcal{K}_A$. Consequently
\[
 1-\lowerdens(G_A)\geq \textup{d}(V_A) .
\]
\end{proposition}
\begin{proof}
Let $N\leq X$ lie in $G_A\cap V_A$. If $N$ is a value of $F_1$ only, then
$N=\sigma(d)$ lies in $\sigma(\N)$, which has density zero by
Lemma \ref{lem:sigma-range-zero}. Otherwise $N=F_e(d)$ with
$2\leq e\leq A$. Put $n=ed\leq AX$ and $k=K_{e,d}\in\mathcal{K}_A$. By
\eqref{eq:fm-congruence}, $N\equiv\sigma(n)\pmod k$, and $k\nmid N$ forces
$k\nmid\sigma(n)$, hence $W_A\nmid\sigma(n)$. The pair $(e,n)$ determines
$N$, so these $N$ number at most
\[
 A\cdot\#\{n\leq AX:W_A\nmid\sigma(n)\}=o_A(X)
\]
by Lemma \ref{lem:fixed-modulus-normality}. Thus all but $o_A(X)$ of the
elements of $V_A\cap[1,X]$ lie outside $G_A$, which gives the density
bound.
\end{proof}
The next estimate determines the scale supplied by this periodic obstruction.
\begin{corollary}\label{cor:fm-envelope-tail}
As $A\to\infty$,
\[
 \Bigl(\frac{e^{-\gamma}}2-o(1)\Bigr)L(A)\leq \textup{d}(V_A)\leq\Delta(A) .
\]
\end{corollary}
\begin{proof}
Every prime $p\leq A$ lies in $\mathcal{K}_A$: with $e=p$ and $d=1$ one has
$D=\{1\}$, $M=C=g=1$ and $K_{p,1}=p$. So $V_A$ consists of $A$-rough
integers, and $\textup{d}(V_A)\leq\Delta(A)$.

For the lower bound let $P$ be the parameter of
Lemma \ref{lem:kernel-tails} with $E=A$. Then
$\log P=(2+o(1))(\log A)(\log\log A)/\log\log\log A$, and Mertens' theorem gives
\[
 \Delta(P)=\Bigl(\frac{e^{-\gamma}}2+o(1)\Bigr)L(A).
\]
It suffices to show that all but a proportion $o(1)$ of the $P$-rough
integers lie in $V_A$.

Let $N$ be $P$-rough and divisible by $k=K_{e,d}$ for some $2\leq e\leq A$.
Then $k$ is $P$-rough. Since $e/g$ divides $k$ and $e/g\leq A<P$, we must
have $e/g=1$, that is, $e\mid M$ and $k=C$. Now $C\geq M\geq e\geq2$, so
$C>P$, and $C\leq(e-1)M$ gives $M>P/A$.

Given $M$, there are at most $\tau(M)$ choices of $e$, since $e\mid M$, and
$e$ and $M$ together determine $D=\{j<e:j\mid M\}$ and hence $C$. So at
most $\tau(M)$ moduli $k$ share a given $M$, and for each of them the
$P$-rough multiples of $C$ have density $\Delta(P)/C\leq\Delta(P)/M$.

Since $M\mid\Lambda(A)$ and $M>P/A$,
\[
 \textup{d}\{N:N\ \text{$P$-rough},\ N\notin V_A\}
 \leq\Delta(P)\sum_{\substack{M\mid\Lambda(A)\\ M>P/A}}\frac{\tau(M)}M
 =\Delta(P)\,S(A,P/A)=o(\Delta(P))
\]
by \eqref{eq:fixed-kernel-tail}. Hence $\textup{d}(V_A)\geq(1-o(1))\Delta(P)$.
\end{proof}

\subsection{Fixed upper levels}

\begin{proof}[Proof of Theorem \ref{thm:almost-log-tail}]
Fix an integer $j\geq3$ and, for sufficiently large $T$, put
\[
 b_j=\frac{j+1}{j-1},\quad
 E=\left\lceil T^{b_j}(\log T)^{3/(j-1)}\right\rceil.
\]
Throughout the counting argument, $j,T,E$ are fixed before $X\to\infty$.
By Proposition \ref{prop:fm-envelope}, all but $o_E(X)$ of the targets in
$V_E\cap[1,X]$ have no representation with cofactor at most $E$.
Lemma \ref{lem:moment} bounds the targets with a representation
$n\leq TN$ and cofactor greater than $E$ by
\[
 C_jT^{j+1}E^{1-j}X\leq\frac{C_jX}{(\log T)^3}.
\]
Moreover, $V_E$ consists of $E$-rough integers. Its nonsquarefree members
up to $X$ therefore number at most
\[
 X\sum_{p>E}\frac1{p^2}\ll\frac XE.
\]
All its members are odd, so Lemma \ref{lem:analytic-odd-representability}
discards only $o(X/\log\log\log X)$ unrepresented targets. Every target surviving
these deletions is represented, squarefree, and satisfies $f(N)>TN$.
Consequently
\begin{align*}
 &\lowerdens\{N\in\Rcal:N\text{ squarefree},\ f(N)>TN\}\\
 &\hspace{15mm}\geq \textup{d}(V_E)-\frac{C_j}{(\log T)^3}-O(E^{-1}).
\end{align*}
Corollary \ref{cor:fm-envelope-tail} and
\[
 L(E)=\left(\frac1{b_j}+o(1)\right)L(T)
 \quad(T\to\infty)
\]
show that the right-hand side is at least
\[
 \left(\frac{e^{-\gamma}}{2b_j}-o(1)\right)
 \frac{\log\log\log T}{(\log T)(\log\log T)}.
\]
Here both error terms are $o(L(T))$ for each fixed $j$. First let
$T\to\infty$ for fixed $j$, and then let $j\to\infty$. Since $b_j\to1$,
this proves \eqref{eq:almost-log-tail}, independently of the finite
representability computation.
\end{proof}

\subsection{Growing upper levels}

\begin{proof}[Proof of Theorem \ref{thm:subexp-growth}]
Fix $0<\varepsilon<1$ sufficiently small in terms of $\eta$. Choose a fixed
$\rho>0$ sufficiently small, depending only on the absolute constants in
Lemma \ref{lem:sigma-rate}, and put
\[
 P=\left\lfloor\rho\frac{\log\log X}{\log\log\log X}\right\rfloor,
 \quad J=\log P,
 \quad W=\sqrt{\frac{J\log J}{2}},
 \quad F=\lfloor e^W\rfloor,
 \quad t=e^{(1-\varepsilon)W}.
\]
In particular $F<P$ for all sufficiently large $X$. Complete periods modulo
$P\#$ and Mertens' theorem give
\begin{equation}\label{eq:sharp-rough-target-count}
 \#\{X/2<N\leq X:\gcd(N,P\#)=1\}
 =\frac X2\Delta(P)+O(P\#)
 \asymp\frac X{\log\log\log X},
\end{equation}
because $P\#=\exp(O(P))=X^{o(1)}$ and $\Delta(P)\asymp1/J$. The
nonsquarefree members of this interval number at most
\[
 X\sum_{p>P}\frac1{p^2}\ll\frac XP=o(X\Delta(P)).
\]

Write $Y=tX$. Lemma \ref{lem:sigma-rate} and the union bound show that the
number of sources $n\leq Y$ for which some prime $p\leq P$ fails to divide
$\sigma(n)$ is at most
\begin{equation}\label{eq:sharp-bad-source-count}
 O(\sqrt Y)+O\left(
 Y\pi(P)\exp\left[-c_1\frac{\log\log Y}{P}\right]
 \right).
\end{equation}
The first term treats $p=2$. The chosen range of $P$ is permitted by
Lemma \ref{lem:sigma-rate} after decreasing $\rho$ if necessary. Since
\[
 Ft=e^{(2-\varepsilon)W}=(\log\log X)^{o(1)},
 \quad
 \exp\left[-c_1\frac{\log\log Y}{P}\right]
 =(\log\log X)^{-c_1/\rho+o(1)},
\]
we may also ensure $c_1/\rho>3$. Multiplying
\eqref{eq:sharp-bad-source-count} by the at most $F$ possible cofactors per
source gives $o(X\Delta(P))$ target--cofactor pairs.

Consider a remaining $P$-rough target $N$ admitting a witness
$N=F_e(n/e)$ with $n\leq tN$ and $e\leq F$. Again $e=1$ is impossible.
For $e\geq2$, put
\[
 L=\lcm\bigl(\{r:r\mid n,\ r<e\}\bigr),\quad
 C=\sum_{\substack{r\mid L\\r<e}}\frac Lr,\quad R=\sigma(n)-N.
\]
The divisors of $L$ below $e$ are exactly the divisors of $n$ below $e$.
By Lemma \ref{lem:fm-modulus},
\[
 R=\frac nL C.
\]
Since
every prime $p\leq P$ divides $\sigma(n)$ but not $N$, the integer
$R=(n/L)C$ is $P$-rough. Writing $u=n/L$, both $u$ and $C$ are therefore
$P$-rough. Since $e\leq F<P$ and $e\mid n=Lu$, every prime-power divisor
of $e$ already divides $L$. Consequently
\[
 n=Lu,\quad \gcd(u,P\#)=1,\quad \gcd(C,P\#)=1,
 \quad e\mid L\mid \Lambda(F).
\]
As before $C\geq L\geq2$, so $C>P$. On the other hand,
\[
 C\leq\sigma(L)
 \leq L\prod_{p\leq F}\left(1-\frac1p\right)^{-1}
 \ll L\log F.
\]
There is therefore an absolute constant $C_0>0$ such that
\begin{equation}\label{eq:moving-kernel-lower-bound}
 L>\frac P{C_0W}.
\end{equation}

Uniformly for $L\mid \Lambda(F)$,
\[
 \log L\leq\log \Lambda(F)\ll F=o(\log X),
 \quad \log(P\#)\ll P=o(\log X).
\]
It follows that $tX/L\geq P\#$ for all sufficiently large $X$. Complete
periods modulo $P\#$ consequently give
\[
 \#\{u\leq tX/L:\gcd(u,P\#)=1\}
 \leq\frac{2tX\Delta(P)}{L}.
\]
For each fixed $L$, the admissible cofactors satisfy $e\mid L$, so there are
at most $\tau(L)$ possibilities. Hence \eqref{eq:moving-kernel-lower-bound}
and Lemma \ref{lem:kernel-tails} bound the entire remaining low-cofactor
witness count by
\begin{align*}
 2tX\Delta(P)
 \sum_{\substack{L\mid \Lambda(F)\\L>P/(C_0W)}}\frac{\tau(L)}L
 &\leq X\Delta(P)\exp\{-\varepsilon W+o(W)\}\\
 &=o(X\Delta(P)).
\end{align*}

Choose a fixed integer $k\geq2$ so large that $(k+1)\varepsilon>2$.
Lemma \ref{lem:moment} bounds the targets with a witness $n\leq tN$ and
cofactor $e>F$ by
\[
 C_kt^{k+1}F^{1-k}X.
\]
After division by $X\Delta(P)$, this is at most
\[
 O_k\left(J\exp\{[2-(k+1)\varepsilon+o(1)]W\}\right)=o(1).
\]
Thus a positive proportion of the squarefree targets in
\eqref{eq:sharp-rough-target-count} has no representation with $n\leq tN$.
These targets are odd. Lemma \ref{lem:analytic-odd-representability}
discards only $o(X/\log\log\log X)=o(X\Delta(P))$ further targets as unrepresented.
Consequently a positive proportion still survives, without any use of the
finite representability computation. Finally,
\[
 \log t
 =\left(\frac{1-\varepsilon}{\sqrt2}+o(1)\right)
 \sqrt{(\log\log\log X)(\log\log\log\log X)}.
\]
Choose $\varepsilon$ sufficiently small in terms of $\eta$. For
$X/2<N\leq X$, replacing $X$ by $N$ in these iterated logarithms changes
the right-hand side by $o(1)$ relative to its main term. This proves
\eqref{eq:subexp-growth}.

For \eqref{eq:positive-moment-growth}, apply the first assertion with a
sufficiently small auxiliary value of $\eta$, raise the resulting ratios to
the power $s$, and note that
\[
 \log\log\log\log X
 =o\left(\sqrt{(\log\log\log X)(\log\log\log\log X)}\right).
\]
The factor $c_\eta/\log\log\log X$ in the proportion of surviving targets can
therefore be absorbed into the prescribed exponent loss $\eta$.
\end{proof}


\section{Bounded representations and cofactor ranges}\label{sec:tightness}

We now study the representations themselves and their contribution to the
distribution of the least representing integer. For $t>0$ and $N\geq1$,
let
\begin{align*}
 r_t(N)&=\#\{(e,d)\in\N^2:F_e(d)=N,\ ed\leq tN\},\\
 r_t^*(N)&=\#\{(e,d)\in\N^2:e\geq2,\ F_e(d)=N,\ ed\leq tN\}.
\end{align*}
The second multiplicity counts proper divisor-prefix sums, excluding the
full divisor sum. Equivalently, $r_t(N)$ counts the pairs $(n,j)$ for
which $N=\sigma_j(n)$ and $n\leq tN$. Thus, on the represented set
$\Rcal$, one has $r_t(N)>0$ if and only if $f(N)\leq tN$.

Recall the empirical measures $\nu_X$ and the normalization $R(X)$ from
\eqref{eq:empirical-measures}. A mass at $\infty$ in a compactified
subsequential limit records escape to infinity. Tightness on the finite
half-line means
\[
 \lim_{T\to\infty}\sup_{X\geq1}\nu_X((T,\infty])=0.
\]

The associated density-one question is
\begingroup
\renewcommand{\theequation}{T}
\begin{equation}\label{eq:T}
 \lim_{A\to\infty}\limsup_{X\to\infty}
 \frac1X\#\{N\leq X:r_A(N)=0\}=0.
\end{equation}
\endgroup
The order of the limits matters: $A$ is fixed before $X\to\infty$.
The next proposition shows that every fixed residue class has many
representation pairs. This does not by itself control the number of
distinct targets.

\subsection{A first moment in every residue class}

For fixed $Q\geq1$ and $\rho\pmod Q$, put
\[
 \underline{\lambda}_{Q,\rho}(A)
 =\liminf_{X\to\infty}\frac QX
 \sum_{\substack{N\leq X\\N\equiv\rho\pmod Q}}r_A(N).
\]

\begin{proposition}
\label{prop:class-first-moment}
For every fixed $Q\geq1$ and every residue class $\rho\pmod Q$,
\[
 \underline{\lambda}_{Q,\rho}(A)
 \gg_{Q,\rho}\frac{A}{(\log A)^2}
 \quad(A\to\infty).
\]
In particular, this lower class mean tends to infinity in every residue
class.
\end{proposition}

\begin{proof}
Fix $Q$ and $\rho$, and set $d_0=\gcd(\rho,Q)$. We first choose a fixed
integer $h$ satisfying
\begin{equation}\label{eq:h-gcd}
 \gcd(\sigma(h),Q)=d_0.
\end{equation}
For each prime power $\ell^k\Vert d_0$, Dirichlet's theorem on primes in
arithmetic progressions permits the choice of a different auxiliary prime
$q_\ell$ congruent to $1$ modulo every prime dividing $Q$. When $\ell=2$,
we also require $q_2\equiv1\pmod4$. Put
\[
 h_\ell=q_\ell^{\ell^k-1},\quad
 h=\prod_{\ell^k\Vert d_0}h_\ell,
\]
with $h=1$ when $d_0=1$. The lifting-the-exponent formula gives
$v_\ell(\sigma(h_\ell))=k$. If $r$ is a prime divisor of $Q$ with $r\ne\ell$, then
$\sigma(h_\ell)\equiv\ell^k\not\equiv0\pmod r$. Multiplicativity of
$\sigma$ proves \eqref{eq:h-gcd}.

The elementary criterion for a linear congruence now gives a unit
$u_\rho\pmod Q$ for which
\begin{equation}\label{eq:u-class}
 \sigma(h)u_\rho\equiv\rho\pmod Q.
\end{equation}
This includes the class $\rho=0$. Prime-power by prime-power, divide by the
common valuation in \eqref{eq:h-gcd}, solve for a unit modulo the remaining
power, and combine the solutions by the Chinese remainder theorem.

Dirichlet's theorem also supplies infinitely many primes $e$ satisfying
\[
 e>\max\{2,h,Q\},\quad e\equiv-1\pmod Q.
\]
For each such prime $e$, define
\[
 \mathcal{S}_e(u_\rho)=\{t\geq1:t\equiv u_\rho\pmod Q,\
 \gcd(t,e\#)=1\}.
\]
The Chinese remainder theorem shows that this periodic set has density
\[
 \delta_e=\frac1Q
 \prod_{\substack{p\leq e\\p\nmid Q}}\left(1-\frac1p\right).
\]
Every prime factor of $t\in\mathcal{S}_e(u_\rho)$ is larger than $e$.

For $t\in\mathcal{S}_e(u_\rho)$, the divisors of $eht$ below $e$ are exactly
the divisors of $h$. Complementary divisors therefore give
\begin{equation}\label{eq:He}
 F_e(ht)=\sigma(h)H_e(t),
 \quad H_e(t)=(e+1)\sigma(t)-et=t+(e+1)s(t),
\end{equation}
where $s(t)=\sigma(t)-t$. Since $Q\mid e+1$, equations
\eqref{eq:u-class} and \eqref{eq:He} give
\[
 F_e(ht)\equiv\sigma(h)t\equiv\rho\pmod Q.
\]
The corresponding representing integer is $n=eht$. Since $H_e(t)\geq t$,
\begin{equation}\label{eq:representing-ratio}
 \frac{n}{F_e(ht)}\leq\frac{eh}{\sigma(h)}\leq eh.
\end{equation}

We next bound the represented value. Complementary divisors give
\[
 \frac{H_e(t)}t
 =1+(e+1)\sum_{\substack{a\mid t\\a>1}}\frac1a.
\]
For a fixed integer $a>1$ with $\gcd(a,e\#)=1$, the Chinese remainder theorem
shows that the set of $t\in\mathcal{S}_e(u_\rho)$ divisible by $a$ has density
$\delta_e/a$. The contribution of a divisor $a$ to the normalized average is at most
$a^{-2}$, since at most $T/a$ integers up to $T$ are divisible by $a$.
Thus $\sum_{a\geq1}a^{-2}<\infty$ justifies interchanging the limit and
the divisor sum.
Therefore
\begin{align*}
 \lim_{T\to\infty}\frac1T
 \sum_{\substack{t\leq T\\t\in\mathcal{S}_e(u_\rho)}}\frac{H_e(t)}t
 &=\delta_e\bigg(1+(e+1)
 \sum_{\substack{a>1\\\gcd(a,e\#)=1}}\frac1{a^2}\bigg)\\
 &\leq3\delta_e.
\end{align*}
Every integer $a>1$ coprime to $e\#$ exceeds $e$. Hence
\[
 (e+1)\sum_{\substack{a>1\\\gcd(a,e\#)=1}}\frac1{a^2}
 \leq(e+1)\sum_{a>e}\frac1{a^2}
 \leq\frac{e+1}{e}\leq\frac32,
\]
which proves the last inequality. Markov's
inequality shows that
\[
 \mathcal{G}_e=\{t\in\mathcal{S}_e(u_\rho):H_e(t)\leq6t\}
\]
has lower density at least $\delta_e/2$.

For a fixed $A$, restrict to primes $e$ satisfying
\[
 \max\{2,h,Q\}<e\leq A/h,
 \quad e\equiv-1\pmod Q.
\]
If $t\in\mathcal{G}_e$ and $t\leq X/(6\sigma(h))$, then the represented value in
\eqref{eq:He} is at most $X$, lies in the class $\rho$, and has a representing integer
$n\leq AN$ by \eqref{eq:representing-ratio}. Distinct pairs $(e,t)$ give
distinct representing integers. Indeed, if $e<e'$ and $eht=e'ht'$, then $e\mid t'$,
contrary to the fact that every prime factor of $t'$ exceeds $e'$. If
$e=e'$, equality of the representing integers forces $t=t'$. Hence these pairs give distinct representation pairs $(n,j)$, even if some
represented values coincide. For each eligible prime $e$, the number of
such $t\leq X/(6\sigma(h))$ is at least
$(\delta_e/2+o(1))X/(6\sigma(h))$. Multiplying by $Q/X$ and summing over
$e$ gives
\begin{equation}\label{eq:class-lower}
 \underline{\lambda}_{Q,\rho}(A)
 \geq\frac1{12\sigma(h)}
 \sum_{\substack{\max\{2,h,Q\}<e\leq A/h\\
 e\text{ prime}\\e\equiv-1\pmod Q}}
 \prod_{\substack{p\leq e\\p\nmid Q}}\left(1-\frac1p\right).
\end{equation}
Mertens' product theorem gives
\[
 \prod_{\substack{p\leq e\\p\nmid Q}}
 \left(1-\frac1p\right)
 \sim\frac Q{\varphi(Q)}\frac{e^{-\gamma}}{\log e}.
\]
For $Z\geq2$, the prime number theorem in arithmetic progressions, followed
by partial summation, gives, as $Z\to\infty$,
\[
 \sum_{\substack{e\leq Z\\e\text{ prime}\\e\equiv-1\pmod Q}}
 \frac1{\log e}
 \sim\frac{Z}{\varphi(Q)(\log Z)^2}.
\]
The omitted primes below the fixed threshold $\max\{2,h,Q\}$ do not
change this asymptotic. Taking $Z=A/h$ in \eqref{eq:class-lower} proves the
proposition.
\end{proof}

\begin{remark}
Proposition \ref{prop:class-first-moment} counts representation pairs. A small set
of represented values could carry many pairs, so the proposition alone gives no bound
for $\#\{N\leq X:r_A(N)=0\}$. Thus the proposition does not by itself imply
\eqref{eq:T}.
\end{remark}

\subsection{Prime cofactor ranges}

The exact forced modulus of Section \ref{sec:arithmetic-preliminaries} always
contains the cofactor when that cofactor is prime. This gives a simple
upper bound for the density of its range.
\begin{corollary}
Let $p$ be prime. Then $F_p(d)\equiv\sigma(pd)\pmod p$ for every $d\geq1$,
and
\[
 \#\bigl(F_p(\N)\cap[1,X]\bigr)\leq\frac Xp+o_p(X),
 \quad\text{so}\quad
 \upperdens\bigl(F_p(\N)\bigr)\leq\frac1p .
\]
\end{corollary}
\begin{proof}
Every $j\in D$ satisfies $j<p$, so $p\nmid M$ and $g=1$. Hence
$p\mid K_{p,d}$, and the congruence follows from \eqref{eq:fm-congruence}.
Now let $N=F_p(d)\leq X$ with $p\nmid N$. Then $p\nmid\sigma(pd)$, and
$pd\leq pN\leq pX$ because $d\leq F_p(d)$. Distinct $N$ have distinct
sources $pd$, so these $N$ number at most
$\#\{n\leq pX:p\nmid\sigma(n)\}$, which is $o_p(X)$ by
Lemma \ref{lem:fixed-modulus-normality}. Every other value of $F_p$ up to
$X$ is a multiple of $p$.
\end{proof}
For $p=2$ the bound $1/2$ is weaker than Proposition \ref{prop:theta-two}
below. For every prime $p$, the numerical upper bound is $1/p$; the error term
$o_p(X)$ is asserted only with $p$ fixed.

\subsection{Bounded cofactors}

For a positive integer $A$, put
\[
 G_A=\bigcup_{e\leq A}F_e(\N),\quad
 \Lambda(A)=\lcm(1,\ldots,A),\quad P_A=A\Lambda(A).
\]
Every $N=F_e(d)$ with
$e\leq A$ has a representing integer $n=ed\leq AN$, because $F_e(d)\geq d$.
Thus
\[
 G_A\subseteq\{N:r_A(N)>0\}.
\]

The upper-tail estimate shows that every fixed union of cofactor ranges misses a positive-density set, even among squarefree targets.

\begin{corollary}
For every fixed $A\geq1$,
\[
 \lowerdens(\N\setminus G_A)>0.
\]
More precisely, as $A\to\infty$,
\begin{equation}\label{eq:fixed-cofactor-defect}
 \lowerdens\{N:N\text{ squarefree},\ N\notin G_A\}
 \geq\left(\frac{e^{-\gamma}}2-o(1)\right)
 \frac{\log\log\log A}{(\log A)(\log\log A)}.
\end{equation}
\end{corollary}

\begin{proof}
If $N\in\Rcal$ and $f(N)>AN$, then $N\notin G_A$. The quantitative
assertion follows from Theorem \ref{thm:almost-log-tail} with $T=A$.
For any fixed smaller $A$, apply that theorem with one sufficiently large
fixed $T\geq A$.
\end{proof}

Write
\[
 \delta_A=\Delta(P_A),\quad
 \eta_A=\upperdens\{N:\gcd(N,P_A\#)>1,\ N\notin G_A\}.
\]
For $2\leq e\leq A$, every forced modulus $K_{e,d}$ has a prime
factor at most $P_A$. Indeed, if $e/g>1$, use a prime factor of $e/g$;
otherwise $K_{e,d}=C$, where $2\leq C\leq(e-1)M\leq P_A$.
There are only finitely many such moduli. Lemma \ref{lem:fm-modulus}
and fixed-modulus normality therefore imply that all but $o_A(X)$
values $F_e(d)\leq X$ have a prime factor at most $P_A$: their sources
satisfy $ed\leq AX$, and a fixed source contributes at most $A$ values.
The case $e=1$ is covered by Lemma \ref{lem:sigma-range-zero}. Hence
\[
 \#\{N\leq X:\gcd(N,P_A\#)=1,\ N\in G_A\}=o_A(X).
\]
Thus $G_A$ has density zero among the integers coprime to
$P_A\#$. Hence the coprime part of $\N\setminus G_A$ has density
$\delta_A$. Since the coprime and noncoprime parts are disjoint,
\begin{equation}\label{eq:bounded-cofactor-complement}
 1-\lowerdens(G_A)=\delta_A+\eta_A.
\end{equation}
Moreover, Mertens' theorem and the prime number theorem, in the form
$\psi(A)=A+o(A)$, give
$\log P_A=\log A+\psi(A)=A+o(A)$ and hence
\begin{equation}\label{eq:delta-A-asymptotic}
 \delta_A=\frac{e^{-\gamma}+o(1)}A.
\end{equation}

The two terms in \eqref{eq:bounded-cofactor-complement} are of different
orders, and it is $\eta_A$ that dominates. Recall the notation
\[
 L(A)=\frac{\log\log\log A}{(\log A)(\log\log A)} .
\]

\begin{proposition}
As $A\to\infty$,
\[
 \eta_A\geq\Bigl(\frac{e^{-\gamma}}2-o(1)\Bigr)L(A),
 \quad\text{so}\quad A\eta_A\to\infty .
\]
\end{proposition}

\begin{proof}
By \eqref{eq:fixed-cofactor-defect},
$1-\lowerdens(G_A)\geq(e^{-\gamma}/2-o(1))L(A)$. Since $AL(A)\to\infty$,
\eqref{eq:delta-A-asymptotic} gives $\delta_A=o(L(A))$. Subtracting in
\eqref{eq:bounded-cofactor-complement} proves the first claim, and the
second follows from $AL(A)\to\infty$.
\end{proof}

One might guess that coprimality to $P_A\#$ is asymptotically
the only obstruction, so that $1-\lowerdens(G_A)\sim e^{-\gamma}/A$. The
proposition rules this out. The integers missed by $G_A$ for other reasons
outnumber the coprime ones by a factor that tends to infinity.

Now set
\[
 \theta_2=\upperdens\bigl(F_2(\N)\bigr).
\]

\begin{proposition}\label{prop:theta-two}
One has
\[
 \theta_2\leq\frac12-0.0602757=0.4397243.
\]
\end{proposition}

\begin{proof}
The only divisor of $2d$ exceeding $d$ is $2d$, and hence
\begin{equation}\label{eq:F2-aliquot}
 F_2(d)=\sigma(2d)-2d=s(2d).
\end{equation}
Thus $F_2(\N)\subseteq s(\N)$ and the cofactor-two range omits every
untouchable number. Write
\[
 \mathcal{U}=\N\setminus s(\N).
\]

We first note that the odd part of $F_2(\N)$ is negligible. If $n$ is even,
then $s(n)$ is odd precisely when $\sigma(n)$ is odd, which occurs precisely
when $n$ is a square or twice a square. For even $n\geq4$ one also has
$s(n)\geq1+n/2$; the exceptional case $n=2$ is harmless. Thus $s(n)\leq X$
implies $n<2X$, apart from at most this one value. Consequently,
\begin{equation}\label{eq:F2-odd-negligible}
 \#\{N\leq X:N\text{ odd},\ N\in F_2(\N)\}=O(\sqrt X).
\end{equation}

Almost all untouchable numbers are even. Indeed, let $N$ be odd. If
$N-1=p+q$ with distinct primes $p$ and $q$, then the proper divisors of $pq$
are $1,p,q$, and therefore $s(pq)=N$. The exceptional-set theorem of
Montgomery and Vaughan shows that the even integers up to $X$ that are not a
sum of two primes number $O(X^{1-\kappa})$ for some $\kappa>0$; the diagonal
possibilities $N-1=2p$ contribute $O(X/\log X)$ more \cite{MVGoldbach}.
Hence
\begin{equation}\label{eq:odd-untouchables}
 \#\{N\leq X:N\text{ odd},\ N\in\mathcal{U}\}=o(X).
\end{equation}
Chen and Zhao proved \cite{ChenZhao} that
\[
 \lowerdens(\mathcal{U})\geq0.0602757.
\]
By \eqref{eq:odd-untouchables}, the number of even
untouchables not exceeding $X$ is at least
$(0.0602757+o(1))X$. There are $X/2+O(1)$ even integers up to $X$, and by
\eqref{eq:F2-odd-negligible} all but $O(\sqrt X)$ elements of
$F_2(\N)\cap[1,X]$ are even. Since \eqref{eq:F2-aliquot} shows that this
range misses every member of $\mathcal{U}$, we obtain
\[
 \#\bigl(F_2(\N)\cap[1,X]\bigr)
 \leq\frac X2-0.0602757X+o(X).
\]
Taking the upper limit proves the proposition.
\end{proof}

\begin{corollary}
For the residual density defined above,
\[
 \eta_2\geq\frac23-0.4397243>0.2269.
\]
\end{corollary}

\begin{proof}
Since $F_1(d)=\sigma(d)$,
\[
 G_2=\sigma(\N)\cup F_2(\N).
\]
Lemma \ref{lem:sigma-range-zero} implies that $\sigma(\N)$ has density
zero, and therefore $\upperdens(G_2)\leq\theta_2$. Moreover,
$\Lambda(2)=2$, $P_2=4$, $P_2\#=6$, and so $\delta_2=1/3$.
Using \eqref{eq:bounded-cofactor-complement},
\[
 \eta_2
 =1-\lowerdens(G_2)-\delta_2
 \geq\frac23-\upperdens(G_2)
 \geq\frac23-0.4397243>0.2269.
\]
\end{proof}

We conjecture that bounded cofactors cover sets of lower density tending to one as the cofactor bound grows.

\begin{conjecture}
\label{conj:bounded-cofactor-weak}
As $A\to\infty$ through the positive integers,
\[
 \lowerdens(G_A)\longrightarrow1.
\]
Equivalently, $\eta_A=o(1)$.
\end{conjecture}

The equivalence in Conjecture \ref{conj:bounded-cofactor-weak} follows from
\eqref{eq:bounded-cofactor-complement} and
\eqref{eq:delta-A-asymptotic}.

Corollary \ref{cor:fm-envelope-tail} shows where the congruence argument
stops. The set $V_A$ consists of $A$-rough integers, so the forced
congruences say nothing about integers with a prime factor at most $A$
that are missed by $G_A$. The $A$-rough part of $\N\setminus G_A$ has
upper density at most $\Delta(A)$, which tends to zero.
Conjecture \ref{conj:bounded-cofactor-weak} is therefore equivalent to the
assertion that the part of $\N\setminus G_A$ with a prime factor at most
$A$ has upper density tending to zero. Every element of that part satisfies some
forced congruence and is nevertheless not a value.

The conjectured density-one coverage by bounded cofactors is equivalent to tightness of the empirical ratio distributions.

\begin{proposition}
\label{prop:tightness-equivalence}
The empirical measures $\nu_X$ form a tight family on $[0,\infty)$
if and only if
\begin{equation}\label{eq:tightness-criterion}
 \lim_{E\to\infty}\upperdens(\N\setminus G_E)=0.
\end{equation}
Equivalently,
\[
 \lim_{E\to\infty}\lowerdens(G_E)=1.
\]
Thus Conjecture \ref{conj:bounded-cofactor-weak} is exactly the assertion
that the family $\{\nu_X\}$ is tight.
\end{proposition}

\begin{proof}
The summand $d$ occurs in $F_e(d)$, so $F_e(d)\geq d$. Hence
\[
 N\in G_E\quad\Longrightarrow\quad f(N)\leq ed\leq EN
 \quad\text{for some }e\leq E.
\]
Using \eqref{eq:exact-representability} to replace the normalization by
$X+O(1)$, we obtain
\[
 \limsup_{X\to\infty}\nu_X((E,\infty])
 \leq\upperdens(\N\setminus G_E).
\]
Thus \eqref{eq:tightness-criterion} makes the tail uniformly small for
all sufficiently large $X$. For $X$ in a bounded interval, only finitely
many represented targets occur, so their ratios have a common finite
upper bound. This proves tightness for the whole family.

Conversely, fix $T\geq1$ and an integer $k\geq2$. Every represented
$N\notin G_E$ with $f(N)\leq TN$ has a witnessing cofactor $e>E$.
Lemma \ref{lem:moment} therefore gives
\[
 \upperdens(\N\setminus G_E)
 \leq\limsup_{X\to\infty}\nu_X((T,\infty])
 +C_kT^{k+1}E^{1-k}.
\]
If the family is tight, first let $E\to\infty$ and then let $T\to\infty$.
The right-hand side tends to zero. Finally,
$\upperdens(\N\setminus G_E)=1-\lowerdens(G_E)$ identically.
\end{proof}

Since $r_A(N)>0$ is equivalent, apart from the two unrepresented integers,
to $f(N)\leq AN$, Proposition \ref{prop:tightness-equivalence} also shows
that the density-one statement \eqref{eq:T} is equivalent to the weak
bounded-cofactor conjecture.


\section{Witness measures and limiting laws}

For a residue class $a\pmod{\Lambda(e)}$, put
\[
 c_e(a)=\sum_{\substack{1\leq j<e\\j\mid a}}\frac1j.
\]
Write $\mathcal{D}$ for the statistical law of
$h(n)=\sigma(n)/n$. The Davenport theorem \cite{Davenport} and its version for arithmetic
progressions by Pollack \cite[Lemma~2]{PollackPalindromes}
gives an atomless measure $\nu_{a,e}$ characterized by
\[
 \nu_{a,e}(B)
 =\lim_{Y\to\infty}\frac1Y
 \#\{n\leq Y:n\equiv a\pmod{\Lambda(e)},\ h(n)\in B\}
\]
for every continuity set $B$. In particular,
each $\nu_{a,e}$ has total mass $1/\Lambda(e)$, and
\begin{equation}
\label{eq:dadd:progression-domination}
 \sum_{a\bmod \Lambda(e)}\nu_{a,e}=\mathcal{D},
 \quad 0\leq\nu_{a,e}\leq\mathcal{D}.
\end{equation}

The progression laws yield explicit continuous limits for the mean numbers of bounded representations.

\begin{proposition}
\label{prop:dadd:witness-means}
For each $t>0$ the limits
\[
 W(t)=\lim_{X\to\infty}\frac1X\sum_{N\leq X}r_t(N),
 \quad
 W^*(t)=\lim_{X\to\infty}\frac1X\sum_{N\leq X}r_t^*(N)
\]
exist, are finite, and are continuous. More explicitly,
\begin{align*}
 W(t)&=\sum_{e\geq1}w_e(t),\quad
 W^*(t)=\sum_{e\geq2}w_e(t), \\
 w_e(t)
 &=\sum_{\substack{a\bmod \Lambda(e)\\e\mid a}}
 \int_{h-c_e(a)\geq1/t}
 \frac{d\nu_{a,e}(h)}{h-c_e(a)}.
\end{align*}
Both series converge uniformly for $0<t\leq T$, for every fixed
$T>0$, and $W(t)\to0$ as $t\downarrow0$.
\end{proposition}

\begin{proof}
For $n\equiv a\pmod{\Lambda(e)}$ with $e\mid a$, complementary divisors
give the exact identities
\[
 F_e(n/e)=n g_e(n),
 \quad
 g_e(n)=h(n)-c_e(a).
\]
In particular $g_e(n)\geq1/e$. A witness with target at most $X$
and ratio at most $t$ is characterized by
\[
 g_e(n)\geq1/t,
 \quad n g_e(n)\leq X.
\]
For completeness, the progression theorem also yields the joint limit
\[
 \frac1X\sum_{\substack{n\leq tX\\n\equiv a\pmod{\Lambda(e)}}}
 \delta_{(n/X,h(n))}
 \ \Longrightarrow\
 \boldsymbol1_{[0,t]}(u)\,du\,d\nu_{a,e}(h).
\]
Indeed, on a rectangle $(\alpha,\beta]\times B$ the assertion follows
by subtracting the progression counts at $\beta X$ and $\alpha X$.
The limiting law in $h$ is tight, so rectangle approximation proves
the joint weak convergence. The boundary of
$\{h-c_e(a)\geq1/t,\ u(h-c_e(a))\leq1\}$ has zero product measure:
the horizontal boundary has zero measure by atomlessness, and for each
$h>c_e(a)$ the curved boundary contains only one value of $u$.
The support condition $h-c_e(a)\geq1/e$ holds for $\nu_{a,e}$ because
it holds for every source in the progression. We may therefore pass to
the limit in the witness condition. Tonelli's theorem gives
\begin{align*}
 &\lim_{X\to\infty}\frac1X
 \#\{n\equiv a\pmod{\Lambda(e)}:
 g_e(n)\geq1/t,\ n g_e(n)\leq X\}\\
 &\quad=
 \int_0^t\nu_{a,e}
 \{h:h-c_e(a)\geq1/t,\ h-c_e(a)\leq1/u\}\,du\\
 &\quad=
 \int_{h-c_e(a)\geq1/t}
 \frac{d\nu_{a,e}(h)}{h-c_e(a)}.
\end{align*}
For the cofactor tail, Lemma \ref{lem:moment}
counts witnessing pairs themselves:
\begin{align*}
 \frac1X\sum_{N\leq X}
 \#\{(e,d):e>E,\ F_e(d)=N,\ ed\leq tN\}
 &\leq\frac{t^k}{X}
 \sum_{n\leq tX}\sum_{\substack{e\mid n\\e>E}}g_e(n)^k
 \nonumber\\
 &\leq C_kt^{k+1}E^{1-k}.
\end{align*}
This proves both limits and uniform convergence on $0<t\leq T$.
Each fixed $w_e(t)$ is continuous by atomlessness, and tends to
zero with $t$ since its integrand is bounded by $t$ on its domain.
The uniform tail proves all remaining assertions.
\end{proof}

Proposition \ref{prop:class-first-moment} with $Q=1$ also shows that
$W(t)\gg t/(\log t)^2$ as $t\to\infty$. Since $w_1(t)\leq1$,
the same lower bound holds for $W^*(t)$.

\subsection{Every subsequential finite law is singular continuous}

All subsequential limits of the empirical ratio distributions have atomless, nonzero finite parts carried by a single Lebesgue-null set.

\begin{theorem}
\label{thm:dadd:universal-singularity}
There is a fixed $\sigma$-compact, Lebesgue-null set
$\mathcal{Z}\subset(0,\infty)$ such that every weak subsequential
limit $\nu$ of $\nu_X$ on $[0,\infty]$ satisfies
\begin{gather*}
 0\in\operatorname{supp}(\nu),\quad \nu(\{0\})=0,\\
 \nu\bigl((0,\infty)\setminus\mathcal{Z}\bigr)=0,\quad
 \nu(\{x\})=0\quad(0<x<\infty).
\end{gather*}
Hence its restriction to $[0,\infty)$ is a nonzero singular
continuous measure. If \eqref{eq:tightness-criterion} holds, every
subsequential limit is a singular continuous probability law on
$[0,\infty)$; in particular, none admits a Lebesgue density.
Neither uniqueness of the limit nor tightness is asserted
unconditionally.
\end{theorem}

\begin{proof}
Erd\H{o}s's singularity theorem asserts that the continuous Davenport
law $\mathcal{D}$ is purely Lebesgue singular; see the account in
\cite[p.~59]{ErdosDistribution}.
By regularity there is a $\sigma$-compact Lebesgue-null set
$K\subset[1,\infty)$ with $\mathcal{D}(K)=1$.
For $c\in\R$, write
\[
 \psi_c(h)=\frac1{h-c}\quad(h>c),
\]
and define the positive witness measure
\[
 \mathcal{W}
 =\sum_{e\geq1}
 \sum_{\substack{a\bmod \Lambda(e)\\e\mid a}}
 (\psi_{c_e(a)})_*
 \bigl(\boldsymbol1_{\{h>c_e(a)\}}
 \psi_{c_e(a)}(h)\,d\nu_{a,e}(h)\bigr).
\]
Proposition \ref{prop:dadd:witness-means} gives
\[
 \mathcal{W}((0,t])=W(t)<\infty\quad(t>0),
\]
so $\mathcal{W}$ is locally finite and atomless. By
\eqref{eq:dadd:progression-domination}, every $\nu_{a,e}$ is
carried by $K$. Therefore $\mathcal{W}$ is carried by
\[
 \mathcal{Z}
 =\bigcup_{e\geq1}
 \bigcup_{\substack{a\bmod \Lambda(e)\\e\mid a}}
 \bigcup_{j\geq1}
 \psi_{c_e(a)}
 \bigl(K\cap[c_e(a)+1/j,j]\bigr).
\]
Each map in the last display is Lipschitz on its indicated set.
Thus $\mathcal{Z}$ is $\sigma$-compact and Lebesgue null.

At finite scale, define the empirical witness measure
\[
 \mathcal{W}_X
 =\frac1X\sum_{N\leq X}
 \sum_{\substack{e,d\geq1\\F_e(d)=N}}
 \delta_{ed/N}.
\]
On every bounded ratio interval it is finite. Proposition
\ref{prop:dadd:witness-means} and continuity of $W$ show that
$\mathcal{W}_X\to\mathcal{W}$ vaguely on $(0,\infty)$. For every
nonnegative $\varphi\in C_c((0,\infty))$, choosing a minimizing
witness separately for each represented target gives
\[
 \int\varphi\,d\nu_X
 \leq\frac X{R(X)}\int\varphi\,d\mathcal{W}_X.
\]
Along any weakly convergent subsequence $\nu_{X_j}\Rightarrow\nu$,
pass to the limit to obtain
\[
 \int\varphi\,d\nu\leq\int\varphi\,d\mathcal{W}.
\]
Hence $\nu|_{(0,\infty)}\leq\mathcal{W}$. This proves both singularity
on the common carrier and the absence of positive finite atoms.

Finally, the same witness domination and portmanteau give, for
every $t>0$,
\[
 \nu(\{0\})
 \leq\nu([0,t))
 \leq\liminf_{j\to\infty}\nu_{X_j}([0,t))
 \leq W(t).
\]
Let $t\downarrow0$ and use
Proposition \ref{prop:dadd:witness-means}. To see that zero nevertheless
belongs to the support, fix $\delta>0$. Theorem \ref{thm:small-values}
and the closed-set direction of portmanteau give
\[
 \nu([0,\delta])
 \geq\limsup_{j\to\infty}\nu_{X_j}([0,\delta])>0.
\]
Applying this with arbitrarily small $\delta$ proves the support claim
and shows that the finite part is nonzero. Finally, tightness excludes
mass at $\infty$ in every compactified subsequential limit, so
Proposition \ref{prop:tightness-equivalence} gives the last assertion.
\end{proof}

The tail estimates imply divergence of positive and logarithmic moments, while inverse moments remain finite and pass to subsequential limits.

\begin{corollary}
Every compactified subsequential law $\nu$ satisfies, as
$T\to\infty$,
\begin{equation}
\label{eq:dadd:subsequential-tail}
 \nu((T,\infty])
 \geq\left(\frac{e^{-\gamma}}2-o(1)\right)
 \frac{\log\log\log T}{(\log T)(\log\log T)}.
\end{equation}
Moreover, for every $s>0$,
\[
 \lim_{X\to\infty}\int x^s\,d\nu_X(x)=\infty,
 \quad
 \int_{[0,\infty]}x^s\,d\nu(x)=\infty.
\]
The logarithmic means also satisfy
\[
 \lim_{X\to\infty}\int\log x\,d\nu_X(x)=\infty,
 \quad
 \int_{[0,\infty]}\log x\,d\nu(x)=\infty.
\]
In contrast, for every $s>0$,
\[
 \sup_{X\geq1}\int x^{-s}\,d\nu_X(x)<\infty,
 \quad
 \int x^{-s}\,d\nu(x)<\infty,
\]
where $\infty^{-s}=0$. Along every weakly convergent subsequence,
the inverse moments converge to those of the limiting law.
In particular, the geometric means of $f(N)/N$ tend to infinity;
even assuming tightness, every subsequential probability law has
infinite logarithmic mean and infinite moments of all positive orders.
\end{corollary}

\begin{proof}
The closed-set direction of portmanteau and
Theorem \ref{thm:almost-log-tail} give
\[
 \nu([T,\infty])
 \geq\limsup_{j\to\infty}\nu_{X_j}([T,\infty])
 \geq\liminf_{X\to\infty}\nu_X((T,\infty]).
\]
Theorem \ref{thm:dadd:universal-singularity} excludes an atom at
$T$, proving \eqref{eq:dadd:subsequential-tail}. For either $\nu_X$
or $\nu$, the $s$th moment is at least $T^s$ times its mass on
$(T,\infty]$. First take the lower limit in $X$ when appropriate,
and then let $T\to\infty$ in the preceding lower bound.

For the logarithmic means, layer cake and Fatou's lemma give
\begin{align*}
 \liminf_{X\to\infty}\int\log^+x\,d\nu_X(x)
 &\geq\int_{T_0}^B
 \liminf_{X\to\infty}\nu_X((T,\infty])\,\frac{dT}{T}\\
 &\geq\left(\frac{e^{-\gamma}}4-o(1)\right)
 (\log\log\log B)^2
 \quad(B\to\infty).
\end{align*}
The same calculation using \eqref{eq:dadd:subsequential-tail}
shows $\int\log^+x\,d\nu(x)=\infty$. On the other hand,
Theorem \ref{thm:small-upper} and the bound $X/R(X)\leq3$ give,
for some $c>0$ and all sufficiently large $u$,
\[
 \sup_X\nu_X([0,e^{-u}])
 \ll\exp\bigl(-\exp(e^{cu})\bigr).
\]
Layer cake therefore bounds $\int\log^-x\,d\nu_X(x)$ uniformly
in $X$. To pass the bound to $\nu$, apply weak convergence to the
continuous bounded functions $\min\{M,\log^-x\}$ on $[0,\infty]$,
defined to equal $M$ at zero and zero at infinity, and then let
$M\to\infty$.
Subtracting the bounded negative parts proves both logarithmic
assertions.

Finally, the same uniform lower-tail bound and layer cake show that
\[
 \sup_{X\geq1}\int_{\{x<1\}}x^{-s}\,d\nu_X(x)
 \leq 1+s\int_0^\infty e^{su}
 \sup_{X\geq1}\nu_X([0,e^{-u}])\,du<\infty.
\]
Apply this also with $2s$ in place of $s$ to obtain uniform integrability
of $x^{-s}$. The bounded truncations $\min\{M,x^{-s}\}$ are continuous
on the compactified half-line when assigned their limiting endpoint
values. Weak convergence, followed by removal of the truncation, proves
the claimed convergence and finiteness of inverse moments.
\end{proof}

\subsection{The collision obstruction}

Convergence of the empirical ratio distributions is equivalent to convergence of the surplus caused by multiple proper witnesses for the same target.

\begin{proposition}
\label{prop:dadd:collision-criterion}
Put
\[
 K_X(t)=\frac1X\sum_{N\leq X}(r_t^*(N)-1)_+.
\]
For every fixed $t>0$,
\begin{equation}
\label{eq:dadd:collision-criterion}
 \nu_X([0,t])=W^*(t)-K_X(t)+o(1).
\end{equation}
Consequently, the threshold density exists if and only if the
genuine proper-witness collision surplus $K_X(t)$ converges.
The whole sequence $\nu_X$ converges weakly on $[0,\infty]$ if and only
if $K_X(t)$ converges for every positive rational $t$.
\end{proposition}

\begin{proof}
A represented target absent from the support of $r_t^*$ can only
be represented at threshold $t$ using cofactor $e=1$; such a
target belongs to $\operatorname{range}(\sigma)$. Lemma
\ref{lem:sigma-range-zero} therefore shows that
\[
 \nu_X([0,t])
 =\frac1X\#\{N\leq X:r_t^*(N)>0\}+o(1).
\]
For every nonnegative integer $r$,
$\boldsymbol1_{\{r>0\}}=r-(r-1)_+$. Sum this identity over
$N\leq X$ with $r=r_t^*(N)$ and apply
Proposition \ref{prop:dadd:witness-means}.

For the final assertion, all compactified subsequential laws have no
finite atoms by Theorem \ref{thm:dadd:universal-singularity}. If the
collision surplus converges at each positive rational threshold, the
displayed identity forces any two subsequential laws to agree on
$[0,t]$ at all those thresholds, hence everywhere, including at infinity.
Compactness then gives convergence of the full sequence. Conversely,
weak convergence implies convergence of the mass of each $[0,t]$,
because its boundary has zero limiting measure, and the identity gives
convergence of $K_X(t)$.
\end{proof}

This collision surplus is bounded away from zero once the ratio threshold is at least two.

\begin{proposition}
For every $t\geq2$,
\[
 \liminf_{X\to\infty}K_X(t)>\frac19.
\]
\end{proposition}

\begin{proof}
Put $Q=5\#=30$, so that $\Delta(5)=\varphi(Q)/Q=4/15$, and let
\[
 \mathcal{C}=\{2,4,8,10\}.
\]
For $D\in\mathcal{C}$ and $\gcd(u,Q)=1$, consider the canonical
cofactor-two witness
\begin{equation}
\label{eq:dadd:collision-affine}
 N=s(Du)=\sigma(D)\sigma(u)-Du
 =u\bigl(s(D)+\sigma(D)(h(u)-1)\bigr).
\end{equation}
Its source $n=Du$ is even, so $N=s(n)\geq n/2$ and its
source-to-target ratio is at most $2\leq t$.

Let $\nu_Q$ be the unnormalized Davenport law restricted to
$\gcd(u,Q)=1$, of total mass $\Delta(5)$. Expanding $h(u)$ as a divisor
sum gives
\[
 \frac1{\Delta(5)}\int(h-1)\,d\nu_Q(h)
 =\sum_{\substack{d\geq1\\\gcd(d,Q)=1}}\frac1{d^2}-1
 =\frac{8\pi^2}{75}-1=:H.
\]
The passage to this first moment is justified by the uniformly bounded
second moment of $h(u)$, which follows from
Lemma \ref{lem:kovac-moment} with $q=3$. The radial
partition argument from Proposition \ref{prop:dadd:witness-means}
gives the source intensity
\begin{align*}
 W_D
 &:=\lim_{X\to\infty}\frac1X
 \#\{u:\gcd(u,Q)=1,\ s(Du)\leq X\}\\
 &=\int\frac{d\nu_Q(h)}{s(D)+\sigma(D)(h-1)}
 \geq\frac{\Delta(5)}{s(D)+\sigma(D)H},
\end{align*}
where the last inequality is Jensen's inequality.

Choose the fixed modulus $q=2^4\cdot3\cdot5^2$.
Lemma \ref{lem:fixed-modulus-normality} implies that, apart
from $o(X)$ sources $n\leq2X$, one has
\[
 v_p(\sigma(n))>v_p(D)=v_p(n)
 \quad(p=2,3,5).
\]
For every remaining target in \eqref{eq:dadd:collision-affine},
\[
 v_p(N)=v_p(\sigma(n)-n)=v_p(D)
 \quad(p=2,3,5).
\]
Thus, up to $o(X)$ exceptions, its image is contained in the
cylinder
\[
 \mathcal{T}_D=\{Dv:\gcd(v,30)=1\},
 \quad \textup{d}(\mathcal{T}_D)=\frac{\Delta(5)}{D}.
\]
If $r_D(N)$ denotes the number of parameters $u$ giving $N$ in
\eqref{eq:dadd:collision-affine}, it follows that
\[
 \liminf_{X\to\infty}\frac1X
 \sum_{N\leq X}(r_D(N)-1)_+
 \geq\Delta(5)\left(
 \frac1{s(D)+\sigma(D)H}-\frac1D\right).
\]
The four cores have distinct valuation vectors at $2,3,5$, whereas
$u$ is coprime to $30$. Thus their source families, and hence their
cofactor-two witnessing pairs, are disjoint. The elementary inequality
\[
 \sum_{D\in\mathcal{C}}(r_D(N)-1)_+
 \leq\left(\sum_{D\in\mathcal{C}}r_D(N)-1\right)_+
 \leq(r_t^*(N)-1)_+
\]
therefore allows the four lower bounds to be added.

Finally, $\pi<22/7$ gives $H<197/3675$. Since
\[
 (D,s(D),\sigma(D))
 \in\{(2,1,3),(4,3,7),(8,7,15),(10,8,18)\},
\]
all four terms are positive and exact rational arithmetic gives
\begin{align*}
 \liminf_{X\to\infty}K_X(t)
 &>
 \frac4{15}
 \sum_{(D,s,B)}
 \left(\frac1{s+197B/3675}-\frac1D\right)\\
 &=\frac{3492203725711}{31001587618275}
 >\frac19.
\end{align*}
This lower bound proves that the collision term in
\eqref{eq:dadd:collision-criterion} is substantial; it does not
prove that the term converges.
\end{proof}


\section{Further questions}

The bounds at the two tails and the structural results above leave several questions about the typical size of the least representing integer.

Theorem \ref{thm:almost-log-tail} gives a fixed-level lower bound of order
\[
 \frac{\log\log\log T}{(\log T)(\log\log T)},
\]
while Theorem \ref{thm:subexp-growth}, for each fixed $\eta>0$, reaches
\[
 \exp\biggl[\left(\frac1{\sqrt2}-\eta\right)
 \sqrt{(\log\log\log N)(\log\log\log\log N)}\biggr]
\]
on
$\gg_\eta X/\log\log\log X$ targets. What are the correct fixed-level tail and the
largest growing scale with a comparable counting estimate? In
particular, does the constant $1/\sqrt2$ reflect the problem itself or the
present estimate for the divisor-kernel tail? At the lower tail, it also
remains to quantify the positive-density result as the ratio tends to zero.

Proposition \ref{prop:tightness-equivalence} identifies tightness with
$\lowerdens(G_E)\to1$. The periodic obstruction $V_E$ lies among the
$E$-rough integers, whose density tends to zero. Tightness is therefore
equivalent to the upper density tending to zero of those integers outside
$G_E$ that have a prime factor at most $E$. Such integers satisfy at least
one of the forced divisibility conditions, but are still not values.
The congruence arguments here do not control this residual set.

The witness means $W(t)$ and $W^*(t)$ exist, and
Proposition \ref{prop:class-first-moment} shows that the first moment of
bounded representations grows in every fixed residue class. Nevertheless,
coverage depends on how these representations are distributed among their
targets. Proposition \ref{prop:dadd:collision-criterion} gives an exact
criterion at each threshold: convergence of the empirical distribution
function is equivalent to convergence of the proper-witness collision
surplus. The lower bound $\liminf K_X(t)>1/9$ for $t\geq2$ confirms that
these collisions cannot be neglected. It provides no convergence result.
A second-moment bound comparable to the square of the first moment could
give positive-density coverage; density-one coverage would require an
additional argument, for example a sufficiently small normalized variance
for a suitable family of witnesses.

Every compactified subsequential law has a singular continuous finite
part on the same null set, with zero in its support and no atom there.
If tightness holds, the resulting probability laws have unbounded support,
infinite logarithmic mean, and infinite moments of every positive order.
It remains to prove or disprove tightness, convergence of the full sequence,
uniqueness of its subsequential limits, and full support on $(0,\infty)$.
The divergence of arithmetic and geometric means is compatible with
tightness and answers none of these existence or uniqueness questions.

We have $\textup{d}(F_1(\N))=0$, $\upperdens(F_2(\N))\leq0.4397243$, and
$\upperdens(F_p(\N))\leq1/p$ for every prime $p$.
Luca and Pomerance's work gives positive lower density for $F_2(\N)$
\cite{LucaPomerance}. It would be useful to decide which other fixed
cofactor ranges have positive lower density, whether their natural
densities exist, and how their overlaps affect the coverage of $G_E$.





\section*{Formalization}



\section*{Methodology and AI usage}
This paper presents the results of a long-term large-scale collaborative project, lasting from November 2025 to August 2026. Its partial results were gradually announced and presented on the forum of the \emph{Erd\H{o}s Problems} website \cite{Bloom1054}, which presents a reliable timeline of project's development. The authors contributed with their thoughts, ideas, proofs, and formalizations, some of which were assisted by generative AI tools. More specifically,
\begin{itemize}
\item E.H. proved the adaptation of the Luca--Pomerance collision sieve to uniformly abundant witnesses in the proof of Theorem \ref{thm:small-values} (Lemmas \ref{lem:sv-regular} and \ref{lem:sv-classes} and Proposition \ref{prop:sv-second-moment}), the exact forced modulus of Lemma \ref{lem:fm-modulus}, and the periodic envelope $V_A$ with its density bounds (Proposition \ref{prop:fm-envelope} and Corollary \ref{cor:fm-envelope-tail}). The proofs of Proposition \ref{prop:class-first-moment}, the complement identity \eqref{eq:bounded-cofactor-complement} with the lower bound for $\eta_A$, Proposition \ref{prop:theta-two} with its corollary, and the prime-cofactor ceiling $\upperdens(F_p(\N))\leq1/p$ were generated with the help of Anthropic's Claude Opus 5 and OpenAI's GPT-5.6 Sol, audited by Claude Fable 5 and GPT-6 Astra, and then verified by E.H. The same models were used for polishing the writeup of E.H.'s parts and for numerical checks.

\item Principia Math's autonomous research harness, using the models GPT-5.5 and
Claude Opus 4.8, resolved the limsup question in Erd\H{o}s Problem~1054,
proving the $\limsup$ variant of Theorem~\ref{thm:almost-log-tail}. The accompanying Lean
formalization includes the circle-method proof of the almost-all
binary Goldbach theorem used as an analytic input. The almost-all Goldbach result was initially useful, but was subsequently made redundant by the proof every integer besides $2$ and $5$ is representable. The formalization of this result, due to H.C. is also included in the Lean formalization repository.
\end{itemize}
The authors independently verified the AI-assisted proofs, carefully inspected the AI-generated parts of the text, rewrote them as appropriate, and they assume full responsibility for the final manuscript. Furthermore, OpenAI's ChatGPT-6 Astra was used for several rounds of copyediting and proofreading.


\section*{Acknowledgments and funding}
We are grateful to Terence Tao for the initial breakthrough that enabled this project and for his encouragement. We would also like to thank Thomas Bloom for maintaining the \emph{Erd\H{o}s Problems} website.

V.K. was supported in part by the Croatian Science Foundation under the project HRZZ-IP-2022-10-5116 (\emph{FANAP}) and in part by the European Union -- NextGenerationEU through the National Recovery and Resilience Plan 2021--2026, via an institutional grant of the University of Zagreb Faculty of Science, IK IA 1.1.3, \emph{Impact4Math}.


\bibliographystyle{plainurl}
\bibliography{EP1054}

\end{document}